\chapter{Feedback Trace and Guarded Trace Semantics}
\label{ch:feedback-trace}

\section{Orientation}
\label{sec:feedback-trace-orientation}

Chapter~\ref{ch:cubical-mealy-execution-concurrency-types} used products to represent
simultaneous execution. Product coordinates expose components executed side by
side. Feedback has a different interface shape. To hide an internal interface,
one separates the visible source and target summands from the hidden summand:
\[
  X+H \longrightarrow Y+H .
\]
The visible source and target interfaces are \(X\) and \(Y\), while \(H\) is
the hidden interface through which execution may circulate internally before
returning to the visible boundary.

Thus this chapter uses the coproduct monoidal structure
\[
  \otimes := +,
  \qquad
  I:=0.
\]
For a morphism
\[
  R:X+H\to Y+H
\]
the feedback operation has the Joyal--Street--Verity type
\[
  \operatorname{Tr}^{H}_{X,Y}:
  \mathcal C(X+H,Y+H)
  \longrightarrow
  \mathcal C(X,Y),
\]
as in the theory of traced monoidal categories
\cite{JoyalStreetVerity1996TracedMonoidalCategories}.

In span semantics, the trace is the finite hidden-path construction
\[
  \operatorname{Tr}^{H}(R)
  =
  R_{XY}
  \boxplus
  R_{HY}\circ
  \mathsf{Path}(R_{HH})
  \circ
  R_{XH}.
\]
Here \(R_{XY}\) records paths that immediately leave through the visible
target. The second summand records paths that enter the hidden interface,
execute zero or more hidden \(H\)-to-\(H\) steps, and then exit through the
visible target.

The block shape is
\[
  R =
  \begin{pmatrix}
    R_{XY} & R_{XH}\\
    R_{HY} & R_{HH}
  \end{pmatrix}
  :
  X+H \longrightarrow Y+H .
\]
The first index records the source summand and the second records the target
summand. Thus \(R_{XH}:X\to H\) is the entrance block, and
\(R_{HY}:H\to Y\) is the exit block. A mnemonic for the four block spans is
\[
\begin{tikzcd}[column sep=huge,row sep=large]
  X_{\mathrm{src}}
    \ar[r,bend left=10,"R_{XY}"]
    \ar[dr,"R_{XH}"']
  &
  Y_{\mathrm{tgt}}
  \\
  H_{\mathrm{src}}
    \ar[ur,"R_{HY}"']
    \ar[r,bend right=10,swap,"R_{HH}"]
  &
  H_{\mathrm{tgt}} .
\end{tikzcd}
\]
This diagram is only a mnemonic for the four restrictions of a single span
over \(X+H\) and \(Y+H\). The source and target decompositions are the
coproduct decompositions
\[
\begin{tikzcd}[column sep=huge,row sep=large]
  X
    \ar[d,phantom,"\oplus" description]
  &
  Y
    \ar[d,phantom,"\oplus" description]
  \\
  H
  &
  H .
\end{tikzcd}
\]
The actual block data are obtained by pulling the source and target legs of
\(R\) back along the corresponding coproduct inclusions.

The guiding distinction is therefore
\[
  \text{product tensor for concurrency,}
  \qquad
  \text{coproduct tensor for feedback.}
\]
Products expose simultaneous coordinates. Coproducts expose visible summands
and allow an internal summand to be hidden by trace. This is the same
graphical intuition underlying string-diagrammatic tensor calculus and
categorical feedback semantics
\cite{JoyalStreet1991GeometryTensorCalculus,KatisSabadiniWalters2002FeedbackTraceFixedPoint},
but in this chapter we use only block and \texttt{tikzcd} diagrams.

This chapter is the unique location where the trace, feedback, finite-path,
guarded-loop, folded-execution, bidirectional signed-interface, and
fixed-point dynamic semantics are developed. Earlier chapters supply the
stateful Mealy, relative execution, and polynomial response-profile
constructions; this chapter supplies the feedback and iteration structure that
acts on those constructions.

The chapter develops the construction in a sequence of layers.

First, we construct the coproduct symmetric monoidal category of spans and its
block-matrix calculus. Then we isolate the finite-path construction
\[
  U\longmapsto \mathsf{Path}(U)
\]
as the free-endospan-monoid construction. The trace on spans is then obtained
by closing a hidden summand and summing over finite hidden paths.

Second, the trace induces the usual iteration operator associated to a
coproduct trace,
\[
  f:X\to Y+X
  \quad\mapsto\quad
  f^\dagger:X\to Y
\]
by solving the recursive \(X\)-summand through finite paths. The connection
between traces and iteration is classical; see
\cite{BloomEsik1993IterationTheories,Hasegawa1997RecursionCyclicSharing,Hasegawa2004Uniformity}.

Third, the coproduct tensor and trace are lifted to internal categories,
Mealy morphisms, and Mealy simulations. The span-level hidden-path normal
forms lift summandwise to Mealy transition-output data and to simulation
cells.

Fourth, the Int construction packages traced feedback into compact closed
signed-interface categories
\[
  (X^+,X^-).
\]
The Int construction is due to Joyal--Street--Verity
\cite{JoyalStreetVerity1996TracedMonoidalCategories}. Its signed-interface
form is closely related to Girard's Geometry of Interaction and its
categorical formulation by Abramsky, Haghverdi, and Scott
\cite{Girard1989GeometryInteractionI,Girard1989TowardsGeometryInteraction,AbramskyHaghverdiScott2002GeometryInteractionLCA}.

Finally, feedback is interpreted logically. Finite paths give the dynamic
logic star, diamonds give least fixed points, boxes give greatest fixed
points, and guarded while programs arise as guarded feedback traces. This is
the dynamic-logic reading of finite program iteration
\cite{HarelKozenTiurynDynamicLogic}, expressed through the subobject-fibration
semantics of categorical logic
\cite{Jacobs1999CategoricalLogic}. Our guarded construction is a
predicate-guarded specialization internal to a total traced span category.
Abstract guarded trace, where trace is available only under guardedness
conditions, is treated in
\cite{GoncharovSchroder2018GuardedTracedCategories}.

We keep three notational levels separate. Composition of spans and Mealy
morphisms is written in ordinary categorical order:
\[
  S\circ R
\]
means first \(R:X\to Y\), then \(S:Y\to Z\). The coproduct tensor is written
\[
  +.
\]
The homwise coproduct of parallel spans is written
\[
  R\boxplus S.
\]
Thus if \(R,S:X\to Y\) are parallel spans, then
\(R\boxplus S:X\to Y\) is the span with apex \(R+S\) and legs induced by the
two pairs of legs.

For objects \(X,Y\), write
\[
  0_{X,Y}
  :=
  \left(
    X\xleftarrow{!}0\xrightarrow{!}Y
  \right)
\]
for the initial span from \(X\) to \(Y\).

The finite-path endospan generated by an endospan \(U:S\to S\) is written
\[
  \mathsf{Path}(U).
\]
The symbols \(;\), \(\cup\), and \((-)^{\star_{\mathrm{path}}}\) are reserved
for the dynamic-logic program syntax introduced later. This notation avoids
conflict with the pullback notation \(f^\ast\).

\section{Feedback span bases}
\label{sec:feedback-span-bases}

The constructions in this chapter use a layered discipline. The purely
feedback-theoretic constructions require finite limits and coproducts
compatible with pullback. The modal interpretation requires additional
regular, Heyting, and \(\omega\)-iterative structure introduced later.

\begin{definition}[Feedback span base]
\label{def:feedback-span-base}
A \textbf{feedback span base} is a finitely complete category \(\mathcal E\)
equipped with the finite and countable coproducts used below such that these
coproducts are disjoint and stable under pullback. Equivalently, whenever
\[
  A\cong\coprod_{i\in I}A_i
\]
is one of the finite or countable coproduct decompositions used in this
chapter, and \(X\to A\) is any morphism, the canonical comparison
\[
  \coprod_{i\in I}(X\times_A A_i)\longrightarrow X
\]
is an isomorphism.

We also require the canonical reindexing isomorphisms for countable
coproducts indexed by countable sets, in particular by \(\mathbb N^2\) and by
the set \(\mathbb N^\ast\) of finite lists of natural numbers. These
reindexing isomorphisms are required to be compatible with pullback-stability
of the coproducts.
\end{definition}

The base condition can be displayed as the following pullback-stable
decomposition:
\[
\begin{tikzcd}[column sep=large,row sep=large]
  X\times_A A_i
    \ar[r]
    \ar[d]
  &
  A_i
    \ar[d,"\iota_i"]
  \\
  X
    \ar[r]
  &
  A .
\end{tikzcd}
\qquad
\leadsto
\qquad
\coprod_i(X\times_A A_i)\cong X .
\]
Consequently pullback distributes over the relevant coproducts. In particular,
for spans
\[
  R:X\to Y,
  \qquad
  S:Y\to Z,
\]
and
\[
  R':X'\to Y',
  \qquad
  S':Y'\to Z',
\]
there is a canonical isomorphism
\[
  (R+R')\times_{Y+Y'}(S+S')
  \cong
  (R\times_Y S)+(R'\times_{Y'}S').
\]
Equivalently,
\[
  (S+S')\circ(R+R')
  \cong
  (S\circ R)+(S'\circ R')
\]
as spans. The comparison can be displayed as
\[
\begin{tikzcd}[column sep=huge,row sep=large]
  (R\times_Y S)+(R'\times_{Y'}S')
    \ar[r,"\cong"]
    \ar[d]
  &
  (R+R')\times_{Y+Y'}(S+S')
    \ar[d]
  \\
  R+R'
    \ar[r,equal]
  &
  R+R' .
\end{tikzcd}
\qquad
\begin{tikzcd}[column sep=huge,row sep=large]
  (R\times_Y S)+(R'\times_{Y'}S')
    \ar[r,"\cong"]
    \ar[d]
  &
  (R+R')\times_{Y+Y'}(S+S')
    \ar[d]
  \\
  S+S'
    \ar[r,equal]
  &
  S+S' .
\end{tikzcd}
\]

All equalities of spans before quotienting by isomorphism are read as
canonical isomorphisms of spans. After passing to isomorphism classes, these
canonical isomorphisms become equalities. Structure maps such as path-flattening are not included in this convention;
they are genuine maps of spans, not generally isomorphisms.

\section{The coproduct symmetric monoidal category of spans}
\label{sec:coproduct-symmetric-monoidal-category-spans}

Let
\[
  \mathbf{Span}^{+}_{\omega}(\mathcal E)
\]
denote the bicategory whose objects are the objects of \(\mathcal E\), whose
1-cells are spans, and whose 2-cells are maps of spans. The underlying
bicategory is the ordinary bicategory of spans. The subscript reminds us that
the countable coproducts used for finite paths are available; the superscript
reminds us that the monoidal structure is coproduct.

The coproduct tensor on objects is
\[
  X\otimes Y:=X+Y,
  \qquad
  I:=0.
\]
For spans
\[
  R=
  \left(
    X\xleftarrow{\ell}R\xrightarrow{r}Y
  \right),
  \qquad
  R'=
  \left(
    X'\xleftarrow{\ell'}R'\xrightarrow{r'}Y'
  \right),
\]
their tensor is the span
\[
  R+R'
  :=
  \left(
    X+X'
    \xleftarrow{\ell+\ell'}
    R+R'
    \xrightarrow{r+r'}
    Y+Y'
  \right).
\]
For maps of spans
\[
  \alpha:R\to S,
  \qquad
  \beta:R'\to S',
\]
their tensor is
\[
  \alpha+\beta:R+R'\to S+S'.
\]

\begin{proposition}[Coproduct tensor as a pseudofunctor]
\label{prop:coproduct-tensor-pseudofunctor-spans}
The coproduct operation extends to a pseudofunctor
\[
  +:
  \mathbf{Span}^{+}_{\omega}(\mathcal E)
  \times
  \mathbf{Span}^{+}_{\omega}(\mathcal E)
  \longrightarrow
  \mathbf{Span}^{+}_{\omega}(\mathcal E).
\]
\end{proposition}

\begin{proof}
\begin{enumerate}
  \item \textbf{Define the object, 1-cell, and 2-cell parts.}
  The object and 1-cell parts are defined by coproduct in \(\mathcal E\). The
  2-cell part is given by coproduct of apex maps.

  \item \textbf{Construct the composition comparison.}
  Let
  \[
    R:X\to Y,\quad S:Y\to Z,\quad
    R':X'\to Y',\quad S':Y'\to Z'
  \]
  be spans. The apex of
  \[
    (S+S')\circ(R+R')
  \]
  is
  \[
    (R+R')\times_{Y+Y'}(S+S').
  \]
  Since the relevant coproducts are disjoint and stable under pullback, the
  canonical comparison
  \[
    (R\times_Y S)+(R'\times_{Y'}S')
    \longrightarrow
    (R+R')\times_{Y+Y'}(S+S')
  \]
  is an isomorphism. Diagrammatically, it is characterized by
  \[
  \begin{tikzcd}[column sep=huge,row sep=large]
    (R\times_Y S)+(R'\times_{Y'}S')
      \ar[r,"\cong"]
      \ar[d]
    &
    (R+R')\times_{Y+Y'}(S+S')
      \ar[d]
    \\
    R+R'
      \ar[r,equal]
    &
    R+R'
  \end{tikzcd}
  \qquad
  \begin{tikzcd}[column sep=huge,row sep=large]
    (R\times_Y S)+(R'\times_{Y'}S')
      \ar[r,"\cong"]
      \ar[d]
    &
    (R+R')\times_{Y+Y'}(S+S')
      \ar[d]
    \\
    S+S'
      \ar[r,equal]
    &
    S+S' .
  \end{tikzcd}
  \]

  \item \textbf{Identify the resulting composite spans.}
  The preceding isomorphism is compatible with source and target legs, so it
  identifies
  \[
    (S+S')\circ(R+R')
    \cong
    (S\circ R)+(S'\circ R').
  \]

  \item \textbf{Construct the unit comparison.}
  The unit comparison is the canonical identification
  \[
    1_{X+X'}\cong 1_X+1_{X'}.
  \]
  In span form this is the isomorphism
  \[
  \begin{tikzcd}[column sep=large,row sep=large]
    X+X'
      \ar[d,equal]
      \ar[r,equal]
    &
    X+X'
      \ar[d,equal]
  \\
    X+X'
      \ar[r,equal]
    &
    X+X' .
  \end{tikzcd}
  \]

  \item \textbf{Verify coherence.}
  The pseudofunctor coherence diagrams compare canonical maps between the same
  finite pullback-coproduct decompositions. They commute by uniqueness of the
  corresponding universal maps.
\end{enumerate}
\end{proof}

\begin{theorem}[Coproduct symmetric monoidal bicategory of spans]
\label{thm:coproduct-symmetric-monoidal-bicategory-spans}
The bicategory
\[
  \mathbf{Span}^{+}_{\omega}(\mathcal E)
\]
is a symmetric monoidal bicategory with tensor \(+\) and unit \(0\).
\end{theorem}

\begin{proof}
\begin{enumerate}
  \item \textbf{Use the tensor pseudofunctor.}
  The tensor pseudofunctor is given by
  Proposition~\ref{prop:coproduct-tensor-pseudofunctor-spans}.

  \item \textbf{Transport coproduct coherence.}
  The associator, left and right unitors, and symmetry are induced by the
  corresponding coproduct isomorphisms in \(\mathcal E\):
  \[
    (X+Y)+Z\cong X+(Y+Z),
    \qquad
    0+X\cong X,
  \]
  \[
    X+0\cong X,
    \qquad
    X+Y\cong Y+X.
  \]

  \item \textbf{View these isomorphisms as invertible spans.}
  Each of these isomorphisms determines an invertible span, and hence an
  equivalence in the bicategory of spans.

  \item \textbf{Check monoidal coherence.}
  The monoidal coherence diagrams are the usual coproduct coherence diagrams
  transported through the span construction. Compatibility with span
  composition is exactly the distributivity comparison of
  Proposition~\ref{prop:coproduct-tensor-pseudofunctor-spans}. Thus the data
  form a symmetric monoidal bicategory.
\end{enumerate}
\end{proof}

Let
\[
  \mathbf{Span}^{+}_{\omega,0}(\mathcal E)
\]
denote the ordinary category with the same objects as
\(\mathbf{Span}^{+}_{\omega}(\mathcal E)\) and with morphisms isomorphism
classes of spans. The preceding theorem descends to a symmetric monoidal
category
\[
  \bigl(\mathbf{Span}^{+}_{\omega,0}(\mathcal E),+,0\bigr).
\]

\section{Block decomposition and matrix calculus}
\label{sec:block-decomposition-feedback}

Let
\[
  R:X+H\to Y+H
\]
be a span
\[
  X+H\xleftarrow{\ell}R\xrightarrow{r}Y+H.
\]
Define four blocks by pulling back the two legs along the coproduct inclusions:
\[
  R_{XY}:=
  X\times_{X+H}R\times_{Y+H}Y,
\]
\[
  R_{XH}:=
  X\times_{X+H}R\times_{Y+H}H,
\]
\[
  R_{HY}:=
  H\times_{X+H}R\times_{Y+H}Y,
\]
and
\[
  R_{HH}:=
  H\times_{X+H}R\times_{Y+H}H.
\]
These are spans
\[
  R_{XY}:X\to Y,
  \qquad
  R_{XH}:X\to H,
\]
\[
  R_{HY}:H\to Y,
  \qquad
  R_{HH}:H\to H.
\]

The four blocks fit into the pullback-stable decomposition
\[
\begin{tikzcd}[column sep=large,row sep=large]
  R_{XY}
    \ar[r]
    \ar[d]
  &
  R
    \ar[d,"{(\ell,r)}"]
  &
  R_{XH}
    \ar[l]
    \ar[d]
  \\
  X\times Y
    \ar[r,"{\iota_X\times\iota_Y}"']
  &
  (X+H)\times(Y+H)
  &
  X\times H
    \ar[l,"{\iota_X\times\iota_H}"]
  \\
  R_{HY}
    \ar[r]
    \ar[d]
  &
  R
    \ar[u,"{(\ell,r)}"']
  &
  R_{HH}
    \ar[l]
    \ar[d]
  \\
  H\times Y
    \ar[r,"{\iota_H\times\iota_Y}"']
  &
  (X+H)\times(Y+H)
  &
  H\times H
    \ar[l,"{\iota_H\times\iota_H}"]
\end{tikzcd}
\]
Equivalently, the apex \(R\) decomposes into the four possible source-target
summand choices.

\begin{proposition}[Block decomposition]
\label{prop:block-decomposition-feedback}
For every span \(R:X+H\to Y+H\), the canonical map
\[
  R_{XY}+R_{XH}+R_{HY}+R_{HH}\longrightarrow R
\]
is an isomorphism of spans over \(X+H\) and \(Y+H\).
\end{proposition}

\begin{proof}
\begin{enumerate}
  \item \textbf{Decompose along the source leg.}
  First decompose \(R\) along its left leg
  \[
    \ell:R\to X+H.
  \]
  The feedback base condition gives
  \[
    R\cong R_X+R_H,
  \]
  where
  \[
    R_X:=X\times_{X+H}R,
    \qquad
    R_H:=H\times_{X+H}R.
  \]

  \item \textbf{Decompose each source summand along the target leg.}
  Now decompose each summand along the right leg into \(Y+H\). Again by
  pullback-stable disjoint coproducts,
  \[
    R_X\cong R_{XY}+R_{XH},
    \qquad
    R_H\cong R_{HY}+R_{HH}.
  \]

  \item \textbf{Assemble the four blocks.}
  Combining these decompositions gives
  \[
    R\cong R_{XY}+R_{XH}+R_{HY}+R_{HH}.
  \]

  \item \textbf{Check compatibility with the span legs.}
  The compatibility with the two span legs follows directly from the pullback
  definitions of the four blocks.
\end{enumerate}
\end{proof}

We write this decomposition as a block matrix
\[
  R=
  \begin{pmatrix}
    R_{XY} & R_{XH}\\
    R_{HY} & R_{HH}
  \end{pmatrix}.
\]
The first index records the source summand and the second records the target
summand.

\begin{lemma}[Matrix multiplication for spans]
\label{lem:span-matrix-multiplication}
Let
\[
  R:X_1+X_2\to Y_1+Y_2,
  \qquad
  S:Y_1+Y_2\to Z_1+Z_2
\]
be spans. Then the \((i,k)\)-block of \(S\circ R\) is canonically
\[
  (S\circ R)_{ik}
  \cong
  \boxplus_{j=1}^{2}
  S_{jk}\circ R_{ij}.
\]
\end{lemma}

\begin{proof}
\begin{enumerate}
  \item \textbf{Start from the composite apex.}
  The apex of \(S\circ R\) is
  \[
    R\times_{Y_1+Y_2}S.
  \]
  Restricting to the \((i,k)\)-block gives the part of this pullback sourced
  at \(X_i\) and targeted at \(Z_k\).

  \item \textbf{Decompose the middle interface.}
  Decompose the \(R\)-factor by its target summand and the \(S\)-factor by its
  source summand:
  \[
    R_i\cong R_{i1}+R_{i2},
    \qquad
    S_k\cong S_{1k}+S_{2k}.
  \]

  \item \textbf{Use disjointness to remove cross terms.}
  Pullback-stable disjoint coproducts give
  \[
    (R_{i1}+R_{i2})\times_{Y_1+Y_2}(S_{1k}+S_{2k})
    \cong
    (R_{i1}\times_{Y_1}S_{1k})
    +
    (R_{i2}\times_{Y_2}S_{2k}).
  \]
  The comparison is displayed by
  \[
  \begin{tikzcd}[column sep=huge,row sep=large]
    (R_{i1}\times_{Y_1}S_{1k})+(R_{i2}\times_{Y_2}S_{2k})
      \ar[r,"\cong"]
      \ar[d]
    &
    (R_i)\times_{Y_1+Y_2}(S_k)
      \ar[d]
    \\
    R_i
      \ar[r,equal]
    &
    R_i .
  \end{tikzcd}
  \qquad
  \begin{tikzcd}[column sep=huge,row sep=large]
    (R_{i1}\times_{Y_1}S_{1k})+(R_{i2}\times_{Y_2}S_{2k})
      \ar[r,"\cong"]
      \ar[d]
    &
    (R_i)\times_{Y_1+Y_2}(S_k)
      \ar[d]
    \\
    S_k
      \ar[r,equal]
    &
    S_k .
  \end{tikzcd}
  \]

  \item \textbf{Read the two summands as span composites.}
  These are precisely the apexes of
  \[
    (S_{1k}\circ R_{i1})
    \boxplus
    (S_{2k}\circ R_{i2}).
  \]
  The source and target maps agree by construction.
\end{enumerate}
\end{proof}

The same proof gives finite block-matrix multiplication for any finite
coproduct decomposition. For two-by-two matrices the rule is the span-valued
analogue of matrix multiplication, with multiplication given by span
composition and addition given by homwise coproduct.

\section{Finite paths and the path monad}
\label{sec:finite-paths-and-path-monad}

Let
\[
  U:S\to S
\]
be an endospan. Define its categorical iterates by
\[
  U^{[0]}:=1_S,
  \qquad
  U^{[n+1]}:=U^{[n]}\circ U.
\]
Thus \(U^{[n]}\) represents \(n\) successive \(U\)-steps.

\begin{definition}[Finite-path endospan]
\label{def:finite-path-endospan}
The \textbf{finite-path endospan} generated by \(U\) is
\[
  \mathsf{Path}(U)
  :=
  \coprod_{n\geq 0}U^{[n]},
\]
where the coproduct is the homwise coproduct of parallel endospans
\(S\to S\).
\end{definition}

This is the linear finite-path/free-monoid analogue of the familiar
polynomial tree and free-monad constructions
\cite{Kock2011PolynomialTrees}; here the generated shapes are finite
linear paths rather than branching trees.

The first few summands are
\[
\begin{tikzcd}[column sep=large,row sep=large]
  S
    \ar[r,"1_S"]
  &
  S
\end{tikzcd}
\qquad
\begin{tikzcd}[column sep=large,row sep=large]
  S
    \ar[r,"U"]
  &
  S
\end{tikzcd}
\qquad
\begin{tikzcd}[column sep=large,row sep=large]
  S
    \ar[r,"U"]
  &
  S
    \ar[r,"U"]
  &
  S .
\end{tikzcd}
\]

\begin{lemma}[Functoriality of finite paths]
\label{lem:path-functoriality-endospans}
Let
\[
  \alpha:U\longrightarrow V
\]
be a map of endospans on \(S\). Define
\[
  \alpha^{[0]}:=1_{1_S},
\]
and recursively
\[
  \alpha^{[n+1]}
  :=
  \alpha^{[n]}\ast_{\mathrm h}\alpha:
  U^{[n+1]}\longrightarrow V^{[n+1]}.
\]
Here \(\ast_{\mathrm h}\) denotes horizontal composition of maps of spans.
Then the coproduct map
\[
  \mathsf{Path}(\alpha)
  :=
  \coprod_{n\geq 0}\alpha^{[n]}
  :
  \mathsf{Path}(U)\longrightarrow\mathsf{Path}(V)
\]
is a map of endospans. This assignment preserves identities and composition.
Hence
\[
  \mathsf{Path}_S:
  \mathrm{EndSpan}_{\mathcal E}(S)
  \longrightarrow
  \mathrm{EndSpan}_{\mathcal E}(S)
\]
is an endofunctor.
\end{lemma}

\begin{proof}
\begin{enumerate}
  \item \textbf{Induce maps on iterated composites.}
  Horizontal composition of span maps is induced by the universal property of
  pullback. Therefore a map
  \[
    \alpha:U\to V
  \]
  induces maps on all iterated composites
  \[
    \alpha^{[n]}:U^{[n]}\to V^{[n]}.
  \]

  \item \textbf{Take the coproduct of the iterated maps.}
  The coproduct universal property gives the map
  \[
    \coprod_{n\geq 0}\alpha^{[n]}:
    \coprod_{n\geq 0}U^{[n]}
    \longrightarrow
    \coprod_{n\geq 0}V^{[n]}.
  \]
  It is a map of endospans because each summand map is a map of endospans.

  \item \textbf{Verify identity preservation.}
  For identities, every iterated map is the identity.

  \item \textbf{Verify composition preservation.}
  For composable maps
  \[
    U\xrightarrow{\alpha}V\xrightarrow{\beta}W,
  \]
  the equality
  \[
    (\beta\alpha)^{[n]}=\beta^{[n]}\alpha^{[n]}
  \]
  follows by induction from functoriality of horizontal composition of span
  maps. Since coproduct injections jointly detect equality, the induced path
  maps preserve identities and composition.
\end{enumerate}
\end{proof}

\begin{lemma}[Two-sided path decomposition]
\label{lem:two-sided-path-decomposition-feedback}
For every endospan \(U:S\to S\), there are canonical isomorphisms
\[
  \mathsf{Path}(U)
  \cong
  1_S\boxplus\bigl(\mathsf{Path}(U)\circ U\bigr)
\]
and
\[
  \mathsf{Path}(U)
  \cong
  1_S\boxplus\bigl(U\circ\mathsf{Path}(U)\bigr).
\]
\end{lemma}

\begin{proof}
\begin{enumerate}
  \item \textbf{Split off the zero-length path for the right decomposition.}
  For the first isomorphism,
  \[
    \mathsf{Path}(U)
    =
    U^{[0]}\boxplus\coprod_{n\geq 0}U^{[n+1]}.
  \]
  Since \(U^{[0]}=1_S\) and \(U^{[n+1]}=U^{[n]}\circ U\), this is
  \[
    1_S
    \boxplus
    \coprod_{n\geq 0}(U^{[n]}\circ U).
  \]

  \item \textbf{Move composition past the coproduct.}
  Composition with \(U\) preserves the countable coproducts used here because
  span composition is pullback and these coproducts are stable under pullback.
  Hence
  \[
    \coprod_{n\geq 0}(U^{[n]}\circ U)
    \cong
    \left(\coprod_{n\geq 0}U^{[n]}\right)\circ U
    =
    \mathsf{Path}(U)\circ U.
  \]

  \item \textbf{Repeat the argument on the other side.}
  For the second isomorphism, use associativity of iterated endospan
  composition to identify
  \[
    U^{[n+1]}\cong U\circ U^{[n]}
  \]
  canonically for each \(n\). The same coproduct argument gives
  \[
    \coprod_{n\geq 0}U^{[n+1]}
    \cong
    U\circ\mathsf{Path}(U).
  \]
\end{enumerate}
\end{proof}

Let
\[
  \mathrm{EndSpan}_{\mathcal E}(S)
  :=
  \mathbf{Span}(\mathcal E)(S,S)
\]
be the hom-category of endospans on \(S\), with tensor product given by span
composition and tensor unit \(1_S\).

\begin{theorem}[Free endospan monoid generated by one step]
\label{thm:free-endospan-monoid-feedback}
The endospan \(\mathsf{Path}(U)\) is the free monoid on \(U\) in the monoidal
category
\[
  \mathrm{EndSpan}_{\mathcal E}(S).
\]
\end{theorem}

\begin{proof}
\begin{enumerate}
  \item \textbf{Define the unit and generator.}
  The monoid unit
  \[
    1_S\longrightarrow\mathsf{Path}(U)
  \]
  is the inclusion of the \(0\)-summand, and the generator
  \[
    U\longrightarrow\mathsf{Path}(U)
  \]
  is the inclusion of the \(1\)-summand.

  \item \textbf{Define multiplication by adding lengths.}
  Since composition distributes over the relevant countable coproducts,
  \[
    \mathsf{Path}(U)\circ\mathsf{Path}(U)
    \cong
    \coprod_{m,n\geq 0}U^{[m]}\circ U^{[n]}.
  \]
  Associativity of span composition gives canonical isomorphisms
  \[
    U^{[m]}\circ U^{[n]}
    \cong
    U^{[m+n]}.
  \]
  The multiplication
  \[
    \mathsf{Path}(U)\circ\mathsf{Path}(U)
    \longrightarrow
    \mathsf{Path}(U)
  \]
  is induced summandwise by these isomorphisms followed by the coproduct
  inclusion of the \((m+n)\)-summand.

  \item \textbf{Construct the universal extension.}
  Let \(M:S\to S\) be a monoid object in
  \(\mathrm{EndSpan}_{\mathcal E}(S)\), with unit
  \[
    \eta:1_S\to M
  \]
  and multiplication
  \[
    \mu:M\circ M\to M.
  \]
  Given a map of endospans
  \[
    f:U\longrightarrow M,
  \]
  define
  \[
    f_n:U^{[n]}\longrightarrow M
  \]
  recursively by
  \[
    f_0:=\eta,
    \qquad
    f_1:=f,
  \]
  and
  \[
    f_{n+1}:
    U^{[n+1]}=U^{[n]}\circ U
    \xrightarrow{f_n\ast_{\mathrm h} f}
    M\circ M
    \xrightarrow{\mu}
    M.
  \]
  The coproduct universal property gives a unique map
  \[
    f^\sharp:\mathsf{Path}(U)\longrightarrow M
  \]
  whose restriction to \(U^{[n]}\) is \(f_n\).

  \item \textbf{Check multiplicativity.}
  We prove by induction on \(n\) that
  \[
    f_{m+n}
    =
    \mu\circ(f_m\ast_{\mathrm h} f_n):
    U^{[m]}\circ U^{[n]}\longrightarrow M
  \]
  after the canonical associativity identification
  \[
    U^{[m]}\circ U^{[n]}\cong U^{[m+n]}.
  \]
  For \(n=0\), this is the right unit law for \(M\). The successor step
  follows from the recursive definition of \(f_{n+1}\) and the associativity
  law for \(\mu\). Hence \(f^\sharp\) preserves multiplication. Preservation
  of the unit is exactly the equation \(f_0=\eta\).

  \item \textbf{Prove uniqueness.}
  Conversely, any monoid morphism
  \[
    g:\mathsf{Path}(U)\to M
  \]
  extending \(f\) is forced to agree with \(\eta\) on the \(0\)-summand and
  with \(f\) on the \(1\)-summand. Multiplicativity then forces the
  restriction to each \(U^{[n]}\) by induction. Thus \(g=f^\sharp\).
\end{enumerate}
\end{proof}

\begin{definition}[Path monad on endospans]
\label{def:path-monad-endospans}
The \textbf{path monad} on endospans over \(S\) is the endofunctor
\[
  \mathsf{Path}_S:
  \mathrm{EndSpan}_{\mathcal E}(S)
  \longrightarrow
  \mathrm{EndSpan}_{\mathcal E}(S)
\]
of Lemma~\ref{lem:path-functoriality-endospans}. Its monad unit
\[
  U\longrightarrow\mathsf{Path}(U)
\]
is the inclusion of the one-step summand, and its monad multiplication
\[
  \mathsf{Path}_S(\mathsf{Path}_S(U))
  \longrightarrow
  \mathsf{Path}_S(U)
\]
is finite-path flattening.
\end{definition}

There are therefore two unit maps in play. The monoid unit of the free monoid
\(\mathsf{Path}(U)\) is the inclusion
\[
  1_S=U^{[0]}\longrightarrow\mathsf{Path}(U),
\]
whereas the monad unit of the free-monoid monad is the inclusion
\[
  U=U^{[1]}\longrightarrow\mathsf{Path}(U).
\]

\begin{proposition}[Path monad laws]
\label{prop:path-monad-laws}
The unit and multiplication of Definition~\ref{def:path-monad-endospans} make
\(\mathsf{Path}_S\) a monad.
\end{proposition}

\begin{proof}
\begin{enumerate}
  \item \textbf{Describe the multiplication summands.}
  The apex of
  \[
    \mathsf{Path}_S(\mathsf{Path}_S(U))
  \]
  decomposes as the coproduct indexed by finite lists of finite
  natural-number lengths. Multiplication concatenates such a list to a single
  finite \(U\)-path using the associativity constraints of span composition.

  \item \textbf{Check associativity on summands.}
  The two sides of the associativity law agree after restriction to each
  summand indexed by a finite list of finite lists of natural numbers. Both
  sides send such a summand to the summand
  \[
    U^{[n_1+\cdots+n_k]}
  \]
  given by the total concatenated length, and the induced map is the same
  iterated associativity isomorphism for span composition.

  \item \textbf{Check the unit laws.}
  The left and right unit laws are the corresponding cases for singleton
  lists and the one-step summand. Since coproduct injections jointly detect
  equality, the monad equations hold.
\end{enumerate}
\end{proof}

\section{The coproduct feedback trace on spans}
\label{sec:coproduct-feedback-trace-spans}

Let
\[
  R:X+H\to Y+H
\]
be a span with block decomposition
\[
  R=
  \begin{pmatrix}
    R_{XY} & R_{XH}\\
    R_{HY} & R_{HH}
  \end{pmatrix}.
\]

\begin{definition}[Coproduct feedback trace]
\label{def:coproduct-feedback-trace}
The \textbf{coproduct feedback trace} of \(R\) over \(H\) is
\[
  \operatorname{Tr}^{H}_{X,Y}(R)
  :=
  R_{XY}
  \boxplus
  R_{HY}\circ\mathsf{Path}(R_{HH})\circ R_{XH}.
\]
Equivalently,
\[
  \operatorname{Tr}^{H}_{X,Y}(R)
  \cong
  R_{XY}
  \boxplus
  \coprod_{n\geq 0}
  R_{HY}\circ R_{HH}^{[n]}\circ R_{XH}.
\]
\end{definition}

The hidden-path summands have the form
\[
\begin{tikzcd}[column sep=large,row sep=large]
  X
    \ar[r,"R_{XH}"]
  &
  H
    \ar[r,"R_{HH}^{[n]}"]
  &
  H
    \ar[r,"R_{HY}"]
  &
  Y .
\end{tikzcd}
\]
The direct summand is
\[
\begin{tikzcd}[column sep=large,row sep=large]
  X
    \ar[r,"R_{XY}"]
  &
  Y .
\end{tikzcd}
\]

\begin{theorem}[Hidden-path normal form]
\label{thm:hidden-path-normal-form}
The span \(\operatorname{Tr}^{H}_{X,Y}(R)\) classifies finite \(R\)-paths
that start in \(X\), end in \(Y\), and whose intermediate interfaces lie in
\(H\).

More explicitly, its apex is canonically the coproduct of:
\begin{enumerate}
  \item the direct visible block \(R_{XY}\);
  \item for each \(n\geq 0\), the composite apex of paths
  \[
    X\xrightarrow{R_{XH}}H
    \xrightarrow{R_{HH}}\cdots
    \xrightarrow{R_{HH}}H
    \xrightarrow{R_{HY}}Y
  \]
  with exactly \(n\) hidden \(H\)-to-\(H\) steps.
\end{enumerate}
In span semantics these paths retain their witnessing apex data; passing to
relational images later forgets this witness multiplicity.
\end{theorem}

\begin{proof}
\begin{enumerate}
  \item \textbf{Identify the direct visible paths.}
  The direct visible paths are represented by the block
  \[
    R_{XY}:X\to Y.
  \]

  \item \textbf{Identify the non-direct hidden paths.}
  A non-direct hidden path first enters \(H\), then takes \(n\) hidden loop
  steps, then exits to \(Y\). Such paths are represented by
  \[
    R_{HY}\circ R_{HH}^{[n]}\circ R_{XH}.
  \]

  \item \textbf{Sum over all finite hidden lengths.}
  Taking the homwise coproduct over \(n\geq 0\) gives
  \[
    R_{HY}\circ\mathsf{Path}(R_{HH})\circ R_{XH}.
  \]

  \item \textbf{Adjoin the direct summand.}
  The homwise coproduct with \(R_{XY}\) therefore classifies exactly the
  finite hidden-path normal forms.
\end{enumerate}
\end{proof}

\begin{lemma}[Finite-word expansion for traced span expressions]
\label{lem:finite-word-expansion-traced-spans}
Every expression built from span composition, coproduct tensor, symmetry, and
a finite sequence of feedback traces admits a canonical expansion, up to the
associator, unitor, symmetry, and coproduct reindexing isomorphisms, as a
coproduct indexed by finite hidden-interface words. For each word, the
associated summand is an iterated span composite determined by the visible
entry, hidden steps, and visible exit appearing in that word.
\end{lemma}

\begin{proof}
\begin{enumerate}
  \item \textbf{Proceed by structural induction.}
  We argue by induction on the construction of the traced expression.

  \item \textbf{Handle composition.}
  Composition concatenates visible composites and pulls back adjacent
  interface summands. The relevant comparison is always a canonical
  associativity isomorphism for iterated pullbacks.

  \item \textbf{Handle coproduct tensor.}
  Coproduct tensor forms disjoint sums of word-indexed families.

  \item \textbf{Handle symmetry.}
  Symmetry renames summands.

  \item \textbf{Handle trace.}
  A trace over a hidden interface replaces a block by the coproduct of all
  finite words in that hidden interface, as in
  Theorem~\ref{thm:hidden-path-normal-form}.

  \item \textbf{Assemble the canonical summandwise maps.}
  The associated summandwise span maps are obtained from associativity of span
  composition, unitors, symmetries, and pullback-coproduct comparison
  isomorphisms. These canonical isomorphisms are the only ambiguity in the
  word-indexed expansion.
\end{enumerate}
\end{proof}

\begin{lemma}[Hidden-path substitution]
\label{lem:hidden-path-substitution}
Let
\[
  R:X+H+K\to Y+H+K
\]
be a span. There is a canonical isomorphism
\[
  \operatorname{Tr}^{H}
  \bigl(
    \operatorname{Tr}^{K}(R)
  \bigr)
  \cong
  \operatorname{Tr}^{H+K}(R).
\]
\end{lemma}

\begin{proof}
\begin{enumerate}
  \item \textbf{Describe the one-step trace.}
  By Theorem~\ref{thm:hidden-path-normal-form}, the right-hand side classifies
  finite \(R\)-paths from \(X\) to \(Y\) whose intermediate interfaces lie in
  \(H+K\).

  \item \textbf{Saturate the \(K\)-interface first.}
  Tracing first over \(K\) replaces every block between the remaining summands
  \(X,H,Y\) by the corresponding \(K\)-saturated block: either a single
  original step not entering \(K\), or a finite path that enters \(K\), follows
  a finite \(K\)-path, and exits \(K\).

  \item \textbf{Trace over \(H\).}
  Tracing the resulting span over \(H\) forms finite strings of these
  \(K\)-saturated \(H\)-steps.

  \item \textbf{Expand the saturated strings.}
  Expanding those finite strings gives exactly the finite paths in the
  original span whose intermediate interfaces lie in \(H+K\).

  \item \textbf{Construct the inverse decomposition.}
  Conversely, every finite path with intermediate interfaces in \(H+K\)
  decomposes uniquely into maximal \(K\)-excursions separated by visits to
  \(H\). This gives inverse canonical maps between the two hidden-path
  coproduct decompositions.

  \item \textbf{Check the span legs.}
  The source and target legs agree under these maps, so the two spans are
  canonically isomorphic.
\end{enumerate}
\end{proof}

\section{The traced span category}
\label{sec:traced-span-category-feedback}

\begin{theorem}[Traced span category]
\label{thm:span-traced-smc}
The symmetric monoidal category
\[
  \bigl(\mathbf{Span}^{+}_{\omega,0}(\mathcal E),+,0\bigr)
\]
is traced, with trace given by
Definition~\ref{def:coproduct-feedback-trace}.
\end{theorem}

\begin{proof}
All equations are equations of isomorphism classes of spans. Before passing
to isomorphism classes, they are canonical isomorphisms.

\begin{enumerate}
  \item \textbf{Verify naturality.}
  Let
  \[
    S:W\to X,
    \qquad
    T:Y\to Z,
    \qquad
    R:X+H\to Y+H.
  \]
  The span
  \[
    (T+1_H)\circ R\circ(S+1_H)
  \]
  has hidden block \(R_{HH}\), entrance block \(R_{XH}\circ S\), exit block
  \(T\circ R_{HY}\), and direct block \(T\circ R_{XY}\circ S\). Thus
  \[
  \begin{aligned}
    \operatorname{Tr}^{H}
    \bigl((T+1_H)\circ R\circ(S+1_H)\bigr)
    &\cong
    T\circ R_{XY}\circ S
    \\
    &\quad\boxplus
    T\circ R_{HY}\circ\mathsf{Path}(R_{HH})
    \circ R_{XH}\circ S
    \\
    &\cong
    T\circ\operatorname{Tr}^{H}(R)\circ S.
  \end{aligned}
  \]
  The equality is represented by the block-composition diagram
  \[
  \begin{tikzcd}[column sep=huge,row sep=large]
    W
      \ar[r,"S"]
    &
    X
      \ar[r,"R_{XH}"]
    &
    H
      \ar[r,"\mathsf{Path}(R_{HH})"]
    &
    H
      \ar[r,"R_{HY}"]
    &
    Y
      \ar[r,"T"]
    &
    Z .
  \end{tikzcd}
  \]

  \item \textbf{Verify dinaturality.}
  Let
  \[
    f:H'\to H
  \]
  be a span and let
  \[
    R:X+H\to Y+H'
  \]
  be a span. The required equation is
  \[
    \operatorname{Tr}^{H}
    \bigl((1_Y+f)\circ R\bigr)
    =
    \operatorname{Tr}^{H'}
    \bigl(R\circ(1_X+f)\bigr).
  \]
  Write the blocks of \(R\) as
  \[
    R_{XY}:X\to Y,\qquad R_{XH'}:X\to H',
  \]
  \[
    R_{HY}:H\to Y,\qquad R_{HH'}:H\to H'.
  \]
  The \(n\)-hidden-step summand on the left is
  \[
    R_{HY}
    \circ
    (f\circ R_{HH'})^{[n]}
    \circ
    f
    \circ
    R_{XH'}.
  \]
  The \(n\)-hidden-step summand on the right is
  \[
    R_{HY}
    \circ
    f
    \circ
    (R_{HH'}\circ f)^{[n]}
    \circ
    R_{XH'}.
  \]
  Associativity of span composition gives, by induction on \(n\),
  \[
    (f\circ R_{HH'})^{[n]}\circ f
    \cong
    f\circ(R_{HH'}\circ f)^{[n]}.
  \]
  The comparison is the finite path-sliding diagram
  \[
  \begin{tikzcd}[column sep=large,row sep=large]
    H'
      \ar[r,"f"]
    &
    H
      \ar[r,"R_{HH'}"]
    &
    H'
      \ar[r,"f"]
    &
    H
      \ar[r,"R_{HH'}"]
    &
    H'
      \ar[r,"\cdots"]
    &
    H'
  \end{tikzcd}
  \]
  rewritten by associativity as
  \[
  \begin{tikzcd}[column sep=large,row sep=large]
    H'
      \ar[r,"f\circ R_{HH'}"]
    &
    H'
      \ar[r,"f\circ R_{HH'}"]
    &
    H'
      \ar[r,"\cdots"]
    &
    H'
  \end{tikzcd}
  \]
  on the left and as
  \[
  \begin{tikzcd}[column sep=large,row sep=large]
    H'
      \ar[r,"f"]
    &
    H
      \ar[r,"R_{HH'}\circ f"]
    &
    H
      \ar[r,"R_{HH'}\circ f"]
    &
    H
      \ar[r,"\cdots"]
    &
    H
  \end{tikzcd}
  \]
  on the right. Thus the \(n\)-summands agree canonically for every \(n\),
  and the direct \(X\to Y\) summand is the same on both sides.

  \item \textbf{Verify Vanishing I.}
  When \(H=0\), all hidden blocks are initial and the trace reduces to the
  visible block \(R\).

  \item \textbf{Verify Vanishing II.}
  Vanishing II is exactly Lemma~\ref{lem:hidden-path-substitution}.

  \item \textbf{Verify superposing.}
  If
  \[
    Q:W\to Z,
  \]
  then \(Q+R\) has \(Q\) entirely in the visible-visible block and has the
  same hidden dynamics as \(R\). Hence
  \[
    \operatorname{Tr}^{H}(Q+R)
    \cong
    Q+\operatorname{Tr}^{H}(R).
  \]

  \item \textbf{Verify yanking.}
  Let
  \[
    \sigma_{H,H}:H+H\to H+H
  \]
  be the symmetry span, with the first copy of \(H\) visible and the second
  copy hidden. Its blocks are
  \[
    R_{\mathrm{vis},\mathrm{vis}}=0_{H,H},
    \qquad
    R_{\mathrm{vis},\mathrm{hid}}=1_H,
  \]
  \[
    R_{\mathrm{hid},\mathrm{vis}}=1_H,
    \qquad
    R_{\mathrm{hid},\mathrm{hid}}=0_{H,H}.
  \]
  Since all positive powers of \(0_{H,H}\) are initial spans,
  \[
    \mathsf{Path}(0_{H,H})
    =
    1_H\boxplus 0_{H,H}\boxplus 0_{H,H}\boxplus\cdots
    \cong 1_H.
  \]
  Therefore the trace is
  \[
    0_{H,H}\boxplus 1_H\circ1_H\circ1_H
    \cong
    1_H.
  \]
  Diagrammatically, the hidden wire is immediately returned:
  \[
  \begin{tikzcd}[column sep=large,row sep=large]
    H
      \ar[r,"1_H"]
    &
    H
      \ar[r,"1_H"]
    &
    H .
  \end{tikzcd}
  \]

  \item \textbf{Conclude the trace axioms.}
  For each Joyal--Street--Verity axiom, expand both sides by
  Lemma~\ref{lem:finite-word-expansion-traced-spans}. The two sides are
  indexed by canonically isomorphic finite hidden-word sets, and the
  corresponding summand spans agree by associativity, unitality, symmetry, and
  pullback-coproduct distributivity. Hence the two sides are canonically
  isomorphic before quotienting and equal after passing to isomorphism
  classes. Thus the trace axioms of
  \cite{JoyalStreetVerity1996TracedMonoidalCategories} hold.
\end{enumerate}
\end{proof}

At the bicategorical level, the same formula defines hom-category trace
functors
\[
  \operatorname{Tr}^{H}_{X,Y}:
  \mathbf{Span}^{+}_{\omega}(\mathcal E)(X+H,Y+H)
  \longrightarrow
  \mathbf{Span}^{+}_{\omega}(\mathcal E)(X,Y).
\]
If
\[
  \alpha:R\to R'
\]
is a map of spans \(R,R':X+H\to Y+H\), then block restriction gives maps
\[
  \alpha_{XY},\alpha_{XH},\alpha_{HY},\alpha_{HH}.
\]
The traced map
\[
  \operatorname{Tr}^{H}(\alpha):
  \operatorname{Tr}^{H}(R)\to \operatorname{Tr}^{H}(R')
\]
is
\[
  \alpha_{XY}
  \boxplus
  \bigl(
    \alpha_{HY}\circ
    \mathsf{Path}(\alpha_{HH})\circ
    \alpha_{XH}
  \bigr),
\]
using functoriality of \(\mathsf{Path}\) and horizontal composition of span
maps. The trace axioms are witnessed by the canonical invertible 2-cells
supplied by the hidden-path normal-form isomorphisms.

\section{The trace-induced Conway dagger}
\label{sec:trace-induced-conway-dagger}

The feedback trace induces an iteration operation for morphisms of the form
\[
  f:X\to Y+X.
\]
Write the two blocks of \(f\) as
\[
  f_Y:X\to Y,
  \qquad
  f_X:X\to X.
\]

\begin{definition}[Trace-induced dagger]
\label{def:trace-induced-dagger-feedback}
The \textbf{trace-induced dagger} of
\[
  f:X\to Y+X
\]
is the span
\[
  f^\dagger
  :=
  f_Y\circ\mathsf{Path}(f_X)
  :
  X\to Y.
\]
\end{definition}

Equivalently, define a span
\[
  L_f:X+X\to Y+X
\]
by
\[
  L_f=
  \begin{pmatrix}
    0_{X,Y} & 1_X\\
    f_Y & f_X
  \end{pmatrix},
\]
where the first copy of \(X\) is visible and the second copy is hidden.

\[
\begin{tikzcd}[column sep=large,row sep=large]
  X_{\mathrm{vis}}
    \ar[r,"1_X"]
  &
  X_{\mathrm{hid}}
    \ar[r,"\mathsf{Path}(f_X)"]
  &
  X_{\mathrm{hid}}
    \ar[r,"f_Y"]
  &
  Y .
\end{tikzcd}
\]

\begin{proposition}[Dagger as trace]
\label{prop:dagger-as-trace-feedback}
There is a canonical isomorphism
\[
  \operatorname{Tr}^{X}(L_f)
  \cong
  f^\dagger.
\]
\end{proposition}

\begin{proof}
\begin{enumerate}
  \item \textbf{Apply the trace formula.}
  By the trace formula,
  \[
    \operatorname{Tr}^{X}(L_f)
    =
    0_{X,Y}
    \boxplus
    f_Y\circ\mathsf{Path}(f_X)\circ1_X.
  \]

  \item \textbf{Remove the initial and identity factors.}
  Removing the initial summand and the identity span gives
  \[
    \operatorname{Tr}^{X}(L_f)
    \cong
    f_Y\circ\mathsf{Path}(f_X)
    =
    f^\dagger.
  \]
\end{enumerate}
\end{proof}

\subsection{Finite approximants and unfolding}
\label{subsec:dagger-approximants-feedback}

The dagger is the collation of all finite approximants.

\begin{definition}[Finite approximants]
\label{def:finite-approximants-dagger}
The \textbf{finite approximants} of \(f^\dagger\) are the spans
\[
  f^{\langle 0\rangle}:=0_{X,Y},
\]
and
\[
  f^{\langle n+1\rangle}
  :=
  f_Y\boxplus
  \bigl(f^{\langle n\rangle}\circ f_X\bigr).
\]
\end{definition}

\begin{lemma}[Approximant expansion]
\label{lem:approximant-expansion-dagger}
For every \(n\geq 0\), there is a canonical isomorphism
\[
  f^{\langle n\rangle}
  \cong
  \coprod_{0\leq k<n}
  f_Y\circ f_X^{[k]}.
\]
Consequently,
\[
  f^\dagger
  \cong
  \coprod_{k\geq 0}
  f_Y\circ f_X^{[k]}.
\]
\end{lemma}

\begin{proof}
\begin{enumerate}
  \item \textbf{Prove the finite expansion by induction.}
  The first statement is by induction on \(n\). For \(n=0\), both sides are
  the initial span \(0_{X,Y}\).

  \item \textbf{Compute the successor case.}
  If the result holds for \(n\), then
  \[
  \begin{aligned}
    f^{\langle n+1\rangle}
    &=
    f_Y\boxplus
    \bigl(f^{\langle n\rangle}\circ f_X\bigr)
    \\
    &\cong
    f_Y
    \boxplus
    \left(
      \coprod_{0\leq k<n}
      f_Y\circ f_X^{[k]}
    \right)
    \circ f_X
    \\
    &\cong
    f_Y
    \boxplus
    \coprod_{0\leq k<n}
    f_Y\circ f_X^{[k+1]}
    \\
    &\cong
    \coprod_{0\leq k<n+1}
    f_Y\circ f_X^{[k]}.
  \end{aligned}
  \]

  \item \textbf{Take the countable collation.}
  The final statement follows from the definition
  \[
    f^\dagger=f_Y\circ\mathsf{Path}(f_X)
  \]
  and the decomposition
  \[
    \mathsf{Path}(f_X)=\coprod_{k\geq 0}f_X^{[k]}.
  \]
\end{enumerate}
\end{proof}

\begin{theorem}[Conway fixpoint identity]
\label{thm:conway-fixpoint-identity-feedback}
For every span
\[
  f:X\to Y+X,
\]
there is a canonical isomorphism
\[
  f^\dagger
  \cong
  f_Y\boxplus(f^\dagger\circ f_X).
\]
Equivalently,
\[
  f^\dagger
  \cong
  [1_Y,f^\dagger]\circ f,
\]
where
\[
  [1_Y,f^\dagger]:Y+X\to Y
\]
is the copairing span obtained from \(1_Y:Y\to Y\) and
\(f^\dagger:X\to Y\).
\end{theorem}

\begin{proof}
\begin{enumerate}
  \item \textbf{Unfold the path object.}
  Using Lemma~\ref{lem:two-sided-path-decomposition-feedback},
  \[
    \mathsf{Path}(f_X)
    \cong
    1_X\boxplus\bigl(\mathsf{Path}(f_X)\circ f_X\bigr).
  \]

  \item \textbf{Compose with the exit block.}
  Composing on the left with \(f_Y\), and using distributivity of composition
  over homwise coproducts, gives
  \[
  \begin{aligned}
    f^\dagger
    &=
    f_Y\circ\mathsf{Path}(f_X)
    \\
    &\cong
    f_Y\circ1_X
    \boxplus
    f_Y\circ\mathsf{Path}(f_X)\circ f_X
    \\
    &\cong
    f_Y\boxplus(f^\dagger\circ f_X).
  \end{aligned}
  \]

  \item \textbf{Rewrite by copairing.}
  The second displayed form is the same equation written after decomposing
  \[
    [1_Y,f^\dagger]\circ f
  \]
  into the \(Y\)-exit and \(X\)-recursive blocks of \(f\).
\end{enumerate}
\end{proof}

\subsection{Naturality, uniformity, and sliding}
\label{subsec:dagger-naturality-uniformity-feedback}

\begin{theorem}[Output naturality]
\label{thm:conway-output-naturality-feedback}
Let
\[
  f:X\to Y+X
\]
and
\[
  g:Y\to Z
\]
be spans. Then
\[
  \bigl((g+1_X)\circ f\bigr)^\dagger
  \cong
  g\circ f^\dagger.
\]
\end{theorem}

\begin{proof}
\begin{enumerate}
  \item \textbf{Identify the recursive and exit blocks.}
  The recursive block of \((g+1_X)\circ f\) is still \(f_X\), while the exit
  block is \(g\circ f_Y\).

  \item \textbf{Apply the dagger formula.}
  Therefore
  \[
    \bigl((g+1_X)\circ f\bigr)^\dagger
    =
    (g\circ f_Y)\circ\mathsf{Path}(f_X)
    \cong
    g\circ f_Y\circ\mathsf{Path}(f_X)
    =
    g\circ f^\dagger.
  \]
\end{enumerate}
\end{proof}

\begin{theorem}[Uniformity]
\label{thm:conway-uniformity-feedback}
Let
\[
  f:X\to Y+X,
  \qquad
  g:X'\to Y+X',
  \qquad
  h:X\to X'
\]
be spans. Suppose there is a canonical isomorphism
\[
  g\circ h
  \cong
  (1_Y+h)\circ f
\]
as spans \(X\to Y+X'\). Then
\[
  g^\dagger\circ h
  \cong
  f^\dagger.
\]
\end{theorem}

\begin{proof}
\begin{enumerate}
  \item \textbf{Decompose the simulation square by blocks.}
  Decomposing the given isomorphism by blocks gives canonical isomorphisms
  \[
    g_Y\circ h\cong f_Y
  \]
  and
  \[
    g_X\circ h\cong h\circ f_X.
  \]

  \item \textbf{Iterate the recursive square.}
  By induction on \(n\),
  \[
    g_X^{[n]}\circ h
    \cong
    h\circ f_X^{[n]}.
  \]
  Diagrammatically, the \(n=1\) square is
  \[
  \begin{tikzcd}[column sep=large,row sep=large]
    X
      \ar[r,"f_X"]
      \ar[d,"h"']
    &
    X
      \ar[d,"h"]
    \\
    X'
      \ar[r,"g_X"']
    &
    X' .
  \end{tikzcd}
  \]

  \item \textbf{Attach the exit block.}
  Thus for every \(n\geq 0\),
  \[
    g_Y\circ g_X^{[n]}\circ h
    \cong
    f_Y\circ f_X^{[n]}.
  \]

  \item \textbf{Sum over all finite lengths.}
  Taking the coproduct over all \(n\geq 0\) gives
  \[
    g_Y\circ\mathsf{Path}(g_X)\circ h
    \cong
    f_Y\circ\mathsf{Path}(f_X),
  \]
  which is precisely
  \[
    g^\dagger\circ h
    \cong
    f^\dagger.
  \]
\end{enumerate}
\end{proof}

\begin{theorem}[Sliding]
\label{thm:conway-sliding-feedback}
Let
\[
  f:X\to Y+Z,
  \qquad
  g:Z\to Y+X
\]
be spans. Let
\[
  \iota_{Y,X}:Y\to Y+X,
  \qquad
  \iota_{Y,Z}:Y\to Y+Z
\]
denote the left coproduct injections. Then
\[
  \bigl([\iota_{Y,X},g]\circ f\bigr)^\dagger
  \cong
  [1_Y,\bigl([\iota_{Y,Z},f]\circ g\bigr)^\dagger]\circ f.
\]
\end{theorem}

\begin{proof}
\begin{enumerate}
  \item \textbf{Expand the left-hand dagger.}
  The left-hand side classifies finite alternating paths that begin with an
  \(f\)-step and then alternate between \(g\)-steps and \(f\)-steps until an
  exit through the \(Y\)-summand occurs.

  \item \textbf{Expand the right-hand dagger.}
  The right-hand side regards the recursive interface as \(Z\) after the
  first \(f\)-step. It therefore classifies the same alternating paths, but
  grouped after the first \(f\)-transition.

  \item \textbf{Compare the word-indexed coproducts.}
  Expanding both daggers by
  Lemma~\ref{lem:approximant-expansion-dagger}, the two sides decompose into
  the same word-indexed coproduct:
  \[
  \begin{tikzcd}[column sep=large,row sep=large]
    X
      \ar[r,"f"]
    &
    Y+Z
      \ar[r,"g\text{ on }Z"]
    &
    Y+X
      \ar[r,"f\text{ on }X"]
    &
    Y+Z
      \ar[r,"\cdots"]
    &
    Y .
  \end{tikzcd}
  \]
  The summands are canonically identified by associativity of span
  composition and the coproduct symmetry.

  \item \textbf{Assemble the canonical isomorphism.}
  The coproduct universal property assembles these summandwise
  identifications into the required canonical isomorphism of spans.
\end{enumerate}
\end{proof}

\subsection{Nested iteration and double dagger}
\label{subsec:dagger-double-feedback}

The double-dagger identity says that two copies of the same recursive
interface may either be solved one after the other or first folded into one
recursive interface and then solved once. In finite-path terms, this is the
unique decomposition of a word in two recursive step types.

\begin{lemma}[Binary path decomposition]
\label{lem:binary-path-decomposition-feedback}
For parallel endospans
\[
  u,v:X\to X,
\]
there is a canonical isomorphism
\[
  \mathsf{Path}(u\boxplus v)
  \cong
  \mathsf{Path}(v)\circ
  \mathsf{Path}\bigl(u\circ\mathsf{Path}(v)\bigr).
\]
It sends a finite word in \(u\) and \(v\) to its unique decomposition into
\(v\)-runs separated by \(u\)-steps.
\end{lemma}

\begin{proof}
\begin{enumerate}
  \item \textbf{Expand the left-hand side as words.}
  The endospan
  \[
    u\boxplus v:X\to X
  \]
  has one-step summands labelled by \(u\) and \(v\). Therefore
  \[
    \mathsf{Path}(u\boxplus v)
  \]
  decomposes as the coproduct over all finite words in the alphabet
  \(\{u,v\}\), with a word interpreted by the corresponding iterated
  composite.

  \item \textbf{Decompose words into \(v\)-runs.}
  With categorical composition read right-to-left operationally, the
  right-hand side classifies words of the form
  \[
    v^{[k_0]}\,u\,v^{[k_1]}\,u\,\cdots\,u\,v^{[k_m]},
  \]
  read operationally from left to right. In categorical composite notation the
  corresponding summand is
  \[
    v^{[k_m]}
    \circ
    (u\circ v^{[k_{m-1}]})
    \circ
    \cdots
    \circ
    (u\circ v^{[k_0]}).
  \]

  \item \textbf{Identify the right-hand side.}
  Thus every finite word in \(u\) and \(v\) is uniquely decomposed into
  \(v\)-runs separated by \(u\)-steps, with a final \(v\)-run. This is
  precisely the summand decomposition of
  \[
    \mathsf{Path}(v)\circ
    \mathsf{Path}\bigl(u\circ\mathsf{Path}(v)\bigr).
  \]

  \item \textbf{Assemble the comparison.}
  On each word-indexed summand the comparison is the canonical associativity
  isomorphism for span composition, and the coproduct universal property
  assembles these summandwise isomorphisms.
\end{enumerate}
\end{proof}

\begin{theorem}[Double dagger]
\label{thm:conway-double-dagger-feedback}
Let
\[
  f:X\to Y+X+X
\]
be a span, and let
\[
  \nabla_X:X+X\to X
\]
be the codiagonal span induced by the codiagonal map in \(\mathcal E\). If
\(f^{\dagger_2}:X\to Y+X\) denotes the span obtained by solving the second
recursive \(X\)-summand first, then there is a canonical isomorphism
\[
  \bigl(f^{\dagger_2}\bigr)^{\dagger_1}
  \cong
  \bigl((1_Y+\nabla_X)\circ f\bigr)^\dagger .
\]
Equivalently, the dagger induced by
Definition~\ref{def:trace-induced-dagger-feedback} satisfies the
double-dagger identity.
\end{theorem}

\begin{proof}
\begin{enumerate}
  \item \textbf{Name the three blocks.}
  Write the three blocks of
  \[
    f:X\to Y+X+X
  \]
  as
  \[
    f_Y:X\to Y,
    \qquad
    f_1:X\to X,
    \qquad
    f_2:X\to X.
  \]

  \item \textbf{Solve the second recursive summand first.}
  Solving the second recursive summand first treats
  \[
    f_Y\boxplus f_1:X\to Y+X
  \]
  as the exit block and \(f_2:X\to X\) as the recursive block. Hence the
  partially solved span
  \[
    f^{\dagger_2}:X\to Y+X
  \]
  has blocks
  \[
    (f^{\dagger_2})_Y
    =
    f_Y\circ\mathsf{Path}(f_2),
    \qquad
    (f^{\dagger_2})_X
    =
    f_1\circ\mathsf{Path}(f_2).
  \]

  \item \textbf{Solve the remaining recursive summand.}
  Solving the remaining recursive summand gives
  \[
    (f^{\dagger_2})^{\dagger_1}
    \cong
    f_Y\circ
    \mathsf{Path}(f_2)
    \circ
    \mathsf{Path}\bigl(f_1\circ\mathsf{Path}(f_2)\bigr).
  \]

  \item \textbf{Solve after folding.}
  Folding the two recursive summands by the codiagonal
  \[
    \nabla_X:X+X\to X
  \]
  gives recursive block
  \[
    f_1\boxplus f_2:X\to X,
  \]
  so
  \[
    \bigl((1_Y+\nabla_X)\circ f\bigr)^\dagger
    \cong
    f_Y\circ\mathsf{Path}(f_1\boxplus f_2).
  \]

  \item \textbf{Apply binary path decomposition.}
  By Lemma~\ref{lem:binary-path-decomposition-feedback}, with \(u=f_1\) and
  \(v=f_2\), there is a canonical isomorphism
  \[
    \mathsf{Path}(f_1\boxplus f_2)
    \cong
    \mathsf{Path}(f_2)
    \circ
    \mathsf{Path}\bigl(f_1\circ\mathsf{Path}(f_2)\bigr).
  \]

  \item \textbf{Attach the exit block.}
  Composing on the left with \(f_Y\) identifies the folded solution with the
  iterated solution. Therefore the double-dagger identity holds.
\end{enumerate}
\end{proof}

In this chapter, a trace-induced Conway operator means a dagger satisfying the
fixpoint, output-naturality, uniformity, sliding, and double-dagger identities
displayed above.

\begin{theorem}[Trace-induced Conway operator]
\label{thm:trace-induced-conway-operator}
The operation
\[
  f\longmapsto f^\dagger
\]
defined by
\[
  f^\dagger=f_Y\circ\mathsf{Path}(f_X)
\]
is the dagger induced by the coproduct trace. It satisfies the Conway
fixpoint, output-naturality, uniformity, sliding, and double-dagger
identities.
\end{theorem}

\begin{proof}
\begin{enumerate}
  \item \textbf{Identify the dagger as a trace.}
  By Proposition~\ref{prop:dagger-as-trace-feedback}, the operation is
  obtained from the coproduct trace.

  \item \textbf{Collect the Conway laws.}
  The fixpoint identity is
  Theorem~\ref{thm:conway-fixpoint-identity-feedback}. Output naturality is
  Theorem~\ref{thm:conway-output-naturality-feedback}. Uniformity is
  Theorem~\ref{thm:conway-uniformity-feedback}. Sliding is
  Theorem~\ref{thm:conway-sliding-feedback}. The double-dagger identity is
  Theorem~\ref{thm:conway-double-dagger-feedback}.

  \item \textbf{Relate the result to the trace/iteration correspondence.}
  This is the span-theoretic instance of the trace/iteration correspondence
  studied in
  \cite{BloomEsik1993IterationTheories,Hasegawa1997RecursionCyclicSharing,Hasegawa2004Uniformity}.
\end{enumerate}
\end{proof}

\section{The coproduct symmetric monoidal bicategory of Mealy morphisms}
\label{sec:coproduct-symmetric-monoidal-bicategory-mealy}

Let \(A\) and \(C\) be internal categories in \(\mathcal E\). Their coproduct
internal category \(A+C\) is defined componentwise:
\[
  (A+C)_0=A_0+C_0,
  \qquad
  (A+C)_1=A_1+C_1.
\]
The identity and composition maps are defined summandwise. The feedback base
condition gives
\[
  (A_1+C_1)\times_{A_0+C_0}(A_1+C_1)
  \cong
  (A_1\times_{A_0}A_1)
  +
  (C_1\times_{C_0}C_1).
\]
Thus there are no cross-summand composable pairs.

\begin{lemma}[No cross arrows in coproduct internal categories]
\label{lem:no-cross-arrows-coproduct-internal-categories}
In the coproduct internal category \(A+C\), every arrow lies entirely in the
\(A\)-summand or entirely in the \(C\)-summand. Equivalently, the mixed
objects
\[
  A_0\times_{A_0+C_0}(A_1+C_1)\times_{A_0+C_0}C_0
\]
and
\[
  C_0\times_{A_0+C_0}(A_1+C_1)\times_{A_0+C_0}A_0
\]
are initial.
\end{lemma}

\begin{proof}
\begin{enumerate}
  \item \textbf{Use the coproduct description of arrows.}
  Since
  \[
    (A+C)_1=A_1+C_1
  \]
  and coproducts are disjoint and stable under pullback, pulling the source
  and target maps back along different summand inclusions gives an initial
  object.

  \item \textbf{Recover the same-summand arrows.}
  Pulling back along the same summand recovers the corresponding arrow object
  of \(A\) or of \(C\). Hence arrows in \(A+C\) have both endpoints in the
  same summand.
\end{enumerate}
\end{proof}

Let
\[
  \mathbf{Mly}^{+}_{\omega}(\mathcal E)
\]
denote the bicategory whose objects are internal categories in \(\mathcal E\),
whose 1-cells are Mealy morphisms, and whose 2-cells are Mealy simulations.
The tensor on objects is componentwise coproduct. For Mealy morphisms
\[
  \Phi:A\rightsquigarrow B,
  \qquad
  \Psi:C\rightsquigarrow D
\]
with carriers
\[
  A_0\xleftarrow{p}P\xrightarrow{q}B_0,
  \qquad
  C_0\xleftarrow{p'}Q\xrightarrow{q'}D_0,
\]
define
\[
  \Phi+\Psi:A+C\rightsquigarrow B+D
\]
to have carrier
\[
  A_0+C_0
  \xleftarrow{p+p'}
  P+Q
  \xrightarrow{q+q'}
  B_0+D_0.
\]
The transition-output structure is defined summandwise.

For simulations
\[
  \Theta:\Phi\Rightarrow\Phi',
  \qquad
  \Lambda:\Psi\Rightarrow\Psi',
\]
define
\[
  \Theta+\Lambda:\Phi+\Psi\Rightarrow\Phi'+\Psi'
\]
by coproduct of the carrier comparison maps and coproduct of the output
comparison arrows.

\begin{theorem}[Coproduct symmetric monoidal bicategory of Mealy morphisms]
\label{thm:coproduct-symmetric-monoidal-bicategory-mealy}
The bicategory
\[
  \mathbf{Mly}^{+}_{\omega}(\mathcal E)
\]
is a symmetric monoidal bicategory with tensor \(+\) and unit the initial
internal category.
\end{theorem}

\begin{proof}
\begin{enumerate}
  \item \textbf{Define the tensor componentwise.}
  The construction on objects, 1-cells, and 2-cells is componentwise
  coproduct.

  \item \textbf{Compare tensor with horizontal composition.}
  By Proposition~\ref{prop:coproduct-tensor-pseudofunctor-spans}, coproduct
  of carrier spans is compatible with horizontal composition up to the
  canonical pullback-coproduct isomorphism
  \[
    (P+Q)\times_{B_0+D_0}(P'+Q')
    \cong
    (P\times_{B_0}P')+(Q\times_{D_0}Q').
  \]
  The comparison is
  \[
  \begin{tikzcd}[column sep=huge,row sep=large]
    (P\times_{B_0}P')+(Q\times_{D_0}Q')
      \ar[r,"\cong"]
      \ar[d]
    &
    (P+Q)\times_{B_0+D_0}(P'+Q')
      \ar[d]
    \\
    P+Q
      \ar[r,equal]
    &
    P+Q .
  \end{tikzcd}
  \]

  \item \textbf{Check the Mealy structure summandwise.}
  The Mealy transition-output maps on the two sides agree summandwise under
  this identification. The unit comparison is inherited from
  \[
    1_{A+C}\cong 1_A+1_C.
  \]

  \item \textbf{Transport the monoidal constraints.}
  The associator, unitors, and symmetry are induced by the coproduct
  isomorphisms of internal categories and carrier spans.

  \item \textbf{Check coherence.}
  The Mealy coherence diagrams reduce to the corresponding coherence diagrams
  for coproducts and to the bicategorical coherence already present for Mealy
  composition. Hence \(\mathbf{Mly}^{+}_{\omega}(\mathcal E)\) is symmetric
  monoidal.
\end{enumerate}
\end{proof}

Let
\[
  \mathbf{Mly}^{+}_{\omega,0}(\mathcal E)
\]
denote the ordinary category whose objects are internal categories in
\(\mathcal E\) and whose morphisms are isomorphism classes of Mealy
morphisms. The preceding theorem descends to a symmetric monoidal category
\[
  \bigl(\mathbf{Mly}^{+}_{\omega,0}(\mathcal E),+,0\bigr).
\]

\section{Coproduct trace for Mealy morphisms}
\label{sec:coproduct-trace-mealy}

Let
\[
  \Phi:A+C\rightsquigarrow B+C
\]
be a Mealy morphism with carrier span
\[
  A_0+C_0
  \xleftarrow{p}
  P
  \xrightarrow{q}
  B_0+C_0.
\]
By block decomposition, the carrier splits into four blocks
\[
  P_{AB},
  \qquad
  P_{AC},
  \qquad
  P_{CB},
  \qquad
  P_{CC}.
\]

\begin{lemma}[Block restriction of Mealy morphisms]
\label{lem:mealy-block-restriction-feedback}
Each carrier block inherits a Mealy morphism:
\[
  \Phi_{AB}:A\rightsquigarrow B,
  \qquad
  \Phi_{AC}:A\rightsquigarrow C,
\]
\[
  \Phi_{CB}:C\rightsquigarrow B,
  \qquad
  \Phi_{CC}:C\rightsquigarrow C.
\]
\end{lemma}

\begin{proof}
\begin{enumerate}
  \item \textbf{Treat one block.}
  We prove the assertion for \(P_{AC}\); the other cases are identical.

  \item \textbf{Restrict the input object.}
  The input object for the restricted transition is
  \[
    (A_1+C_1)\times_{A_0+C_0}P_{AC}.
  \]
  Since \(P_{AC}\) lies over the \(A_0\)-summand on the left,
  pullback-stable disjoint coproducts give
  \[
    (A_1+C_1)\times_{A_0+C_0}P_{AC}
    \cong
    A_1\times_{A_0}P_{AC}.
  \]

  \item \textbf{Show the transition stays in the source block.}
  The Mealy transition equation for the input interface says that the new left
  interface is obtained from the endpoint of the input arrow. Since the input
  arrow lies in \(A_1\), the new left interface lies in \(A_0\).

  \item \textbf{Use absence of cross arrows for the output.}
  The old output object \(q(x)\) lies in \(C_0\). The Mealy output arrow has
  the standing orientation
  \[
    \omega(a,x):q(\delta(a,x))\longrightarrow q(x).
  \]
  Thus it is an arrow of \(B+C\) whose target lies in \(C_0\). By
  Lemma~\ref{lem:no-cross-arrows-coproduct-internal-categories}, there are no
  arrows from the \(B\)-summand to the \(C\)-summand or conversely. Hence
  \(q(\delta(a,x))\) also lies in \(C_0\), and the output arrow lies in
  \(C_1\).

  \item \textbf{Express the restricted structure categorically.}
  Under the decomposition
  \[
    (A_1+C_1)\times_{A_0+C_0}P_{AC}
    \cong
    A_1\times_{A_0}P_{AC},
  \]
  the original Mealy transition-output map factors through the corresponding
  restricted transition-output object for a Mealy morphism
  \(A\rightsquigarrow C\). This factorization is unique because block
  inclusions are disjoint and jointly detect maps out of the decomposed
  domain.

  \item \textbf{Restrict the equations.}
  The unit and multiplication axioms restrict because block inclusions jointly
  detect equality after the block decomposition.
\end{enumerate}
\end{proof}

\begin{lemma}[Countable homwise sums of parallel Mealy morphisms]
\label{lem:homwise-countable-coproducts-mealy}
Let
\[
  (\Phi_n)_{n\geq 0}:A\rightsquigarrow B
\]
be a countable family of parallel Mealy morphisms. Under the feedback span
base hypotheses, the carrier coproduct
\[
  \coprod_{n\geq 0}P_n
\]
carries a canonical Mealy morphism
\[
  \boxplus_{n\geq 0}\Phi_n:A\rightsquigarrow B
\]
whose transition and output maps restrict to those of \(\Phi_n\) on the
\(n\)-th summand. This construction is functorial in Mealy simulations and is
preserved by horizontal composition in each variable up to the canonical
pullback-coproduct isomorphisms. It is a homwise sum construction; no
universal property in the simulation hom-category is asserted here.
\end{lemma}

\begin{proof}
\begin{enumerate}
  \item \textbf{Decompose the transition object.}
  The transition object for the carrier coproduct is
  \[
    A_1\times_{A_0}\left(\coprod_{n\geq 0}P_n\right).
  \]
  By pullback-stability of the relevant countable coproducts, this is
  canonically isomorphic to
  \[
    \coprod_{n\geq 0}(A_1\times_{A_0}P_n).
  \]

  \item \textbf{Define the structure summandwise.}
  Define the transition-output map summandwise using the corresponding
  structure map of \(\Phi_n\).

  \item \textbf{Verify the Mealy equations.}
  The Mealy unit and multiplication equations hold after restriction to every
  summand, hence globally because coproduct injections jointly detect
  equality. The same argument applies to simulations.

  \item \textbf{Check compatibility with horizontal composition.}
  Compatibility with horizontal composition follows because the carrier of a
  horizontal composite is a pullback, and pullback distributes over the
  countable coproducts used here.
\end{enumerate}
\end{proof}

For an endo-Mealy morphism
\[
  \Psi:C\rightsquigarrow C,
\]
define
\[
  \Psi^{[0]}:=1_C,
  \qquad
  \Psi^{[n+1]}:=\Psi^{[n]}\circ\Psi,
\]
and
\[
  \mathsf{Path}(\Psi):=\boxplus_{n\geq 0}\Psi^{[n]}.
\]
This is well-defined by
Lemma~\ref{lem:homwise-countable-coproducts-mealy}.

\begin{definition}[Coproduct trace of a Mealy morphism]
\label{def:mealy-coproduct-trace}
For
\[
  \Phi:A+C\rightsquigarrow B+C,
\]
the \textbf{coproduct trace} of \(\Phi\) over \(C\) is
\[
  \operatorname{Tr}^{C}_{A,B}(\Phi)
  :=
  \Phi_{AB}
  \boxplus
  \Phi_{CB}\circ\mathsf{Path}(\Phi_{CC})\circ\Phi_{AC}.
\]
\end{definition}

\begin{theorem}[Traced Mealy category]
\label{thm:mly-traced-smc}
The symmetric monoidal category
\[
  \bigl(\mathbf{Mly}^{+}_{\omega,0}(\mathcal E),+,0\bigr)
\]
is traced. The trace is given by
Definition~\ref{def:mealy-coproduct-trace}.
\end{theorem}

\begin{proof}
\begin{enumerate}
  \item \textbf{Restrict to blocks.}
  By Lemma~\ref{lem:mealy-block-restriction-feedback}, every Mealy morphism
  \[
    \Phi:A+C\rightsquigarrow B+C
  \]
  decomposes into four Mealy blocks.

  \item \textbf{Build the hidden-path Mealy morphism.}
  For each \(n\geq 0\), the composite
  \[
    \Phi_{CB}\circ\Phi_{CC}^{[n]}\circ\Phi_{AC}
  \]
  is a Mealy morphism \(A\rightsquigarrow B\). The homwise sum over all
  \(n\), together with the direct block \(\Phi_{AB}\), is again a Mealy
  morphism by Lemma~\ref{lem:homwise-countable-coproducts-mealy}.

  \item \textbf{Check the Mealy laws summandwise.}
  The Mealy unit and multiplication laws hold summandwise. On each hidden-path
  summand, the transition-output structure is the iterated composite
  transition-output structure of the corresponding finite Mealy path. Since
  coproduct injections jointly detect equality, the laws hold for the whole
  trace.

  \item \textbf{Lift the span trace axioms.}
  The trace axioms are inherited from the span trace axioms by hidden-path
  normal forms. For each trace axiom, the two sides decompose into the same
  finite hidden Mealy paths. The carrier spans are canonically isomorphic by
  Theorem~\ref{thm:span-traced-smc}, and under these isomorphisms the Mealy
  transition maps agree summandwise.

  \item \textbf{Check output data.}
  The output data agree summandwise because both sides apply the same iterated
  Mealy transition-output composite along the same hidden path: the
  intermediate \(C\)-outputs are consumed as the inputs of the next hidden
  block, and the final \(B\)-output is produced by the same exit block.

  \item \textbf{Pass to isomorphism classes.}
  Therefore the canonical span isomorphisms lift to isomorphisms of Mealy
  morphisms. Passing to isomorphism classes gives the Joyal--Street--Verity
  trace axioms.
\end{enumerate}
\end{proof}

\section{Trace on Mealy simulations}
\label{sec:trace-simulations-feedback}

Let
\[
  \Theta:\Phi\Rightarrow\Psi
\]
be a Mealy simulation between Mealy morphisms
\[
  \Phi,\Psi:A+C\rightsquigarrow B+C.
\]
Let \(P\) be the carrier of \(\Phi\), and let \(Q\) be the carrier of
\(\Psi\). We write the output-comparison Kleisli datum in component form as a
carrier comparison
\[
  \theta:Q\to P
\]
together with an output-comparison arrow
\[
  \alpha:Q\to (B+C)_1.
\]
These satisfy the source, target, input-processing, and output-lax equations
of the Mealy simulation definition:
\[
  p_P\theta=p_Q,
  \qquad
  s_{B+C}\alpha=q_P\theta,
  \qquad
  t_{B+C}\alpha=q_Q,
\]
together with the transition-output compatibility equations.

\begin{lemma}[Block restriction of simulations]
\label{lem:simulation-block-restriction-feedback}
The simulation \(\Theta\) decomposes into block simulations
\[
  \Theta_{AB},
  \qquad
  \Theta_{AC},
  \qquad
  \Theta_{CB},
  \qquad
  \Theta_{CC}.
\]
\end{lemma}

\begin{proof}
\begin{enumerate}
  \item \textbf{Treat one block.}
  We prove the assertion for the \(AC\)-block.

  \item \textbf{Use preservation of the input interface.}
  Let \(y\in Q_{AC}\). Then the left interface of \(y\) lies in \(A_0\), and
  its output interface lies in \(C_0\). The carrier comparison \(\theta\)
  preserves the left interface, so \(\theta(y)\) also lies over \(A_0\) on
  the left.

  \item \textbf{Use absence of cross arrows for the output comparison.}
  The output-comparison arrow \(\alpha_y\) has endpoints the output interface
  of \(\theta(y)\) and the output interface of \(y\), both regarded as objects
  of \(B+C\). Since the output interface of \(y\) lies in \(C_0\), and \(B+C\)
  has no cross-summand arrows, the output interface of \(\theta(y)\) also
  lies in \(C_0\). Hence \(\theta(y)\in P_{AC}\), and \(\alpha_y\) lies in
  \(C_1\).

  \item \textbf{Restrict the simulation equations.}
  The source, target, input-processing, and output-lax equations restrict
  because block inclusions jointly detect equality after the coproduct
  decomposition. The same argument applies to the other blocks.
\end{enumerate}
\end{proof}

\begin{lemma}[Coproducts and composites of simulations]
\label{lem:coproducts-composites-simulations-feedback}
Mealy simulations are closed under horizontal composition, finite and
countable homwise sums, and block restriction along coproduct decompositions.
Moreover these operations are compatible with vertical composition.
\end{lemma}

\begin{proof}
\begin{enumerate}
  \item \textbf{Use horizontal composition in the Mealy bicategory.}
  Horizontal composition is the horizontal composition of Mealy simulations
  from the bicategory of Mealy morphisms.

  \item \textbf{Form homwise sums summandwise.}
  Homwise sums are formed by taking coproducts of carrier comparison maps and
  coproducts of output-comparison arrows. Since the relevant pullbacks
  distribute over these coproducts, the simulation equations hold summandwise
  and hence globally.

  \item \textbf{Use block restriction.}
  Block restriction is Lemma~\ref{lem:simulation-block-restriction-feedback}.

  \item \textbf{Check compatibility with vertical composition.}
  Compatibility with vertical composition follows because vertical composition
  is ordinary composition of the carrier comparison maps together with
  composition of the output-comparison arrows, and coproduct decomposition is
  obtained by pullback along coproduct injections. Pullback commutes with
  these compositions.
\end{enumerate}
\end{proof}

For each \(n\geq 0\), define
\[
  \Theta_{CC}^{[n]}:
  \Phi_{CC}^{[n]}\Rightarrow\Psi_{CC}^{[n]}
\]
by iterated horizontal composition of simulations, with
\[
  \Theta_{CC}^{[0]}:=1_{1_C}.
\]

\begin{lemma}[Finite hidden composites of simulations]
\label{lem:finite-hidden-composites-simulations}
For every \(n\geq 0\), the horizontal composite
\[
  \Theta_{CB}
  \ast_{\mathrm h}
  \Theta_{CC}^{[n]}
  \ast_{\mathrm h}
  \Theta_{AC}
\]
is a Mealy simulation
\[
  \Phi_{CB}\circ\Phi_{CC}^{[n]}\circ\Phi_{AC}
  \Rightarrow
  \Psi_{CB}\circ\Psi_{CC}^{[n]}\circ\Psi_{AC}.
\]
These composites are compatible with vertical composition by interchange.
\end{lemma}

\begin{proof}
\begin{enumerate}
  \item \textbf{Use block simulations.}
  Each displayed block simulation exists by
  Lemma~\ref{lem:simulation-block-restriction-feedback}.

  \item \textbf{Compose horizontally.}
  Horizontal composition of Mealy simulations supplies the displayed composite
  simulation. For \(n=0\), the middle factor is the identity simulation on the
  identity Mealy morphism. For \(n+1\), the middle factor is obtained by
  composing \(\Theta_{CC}^{[n]}\) with \(\Theta_{CC}\).

  \item \textbf{Use interchange.}
  Compatibility with vertical composition is the interchange law in the
  bicategory of Mealy morphisms. The source, target, input-processing, and
  output-lax equations are preserved by horizontal composition.
\end{enumerate}
\end{proof}

Because the displayed simulation direction is opposite to the carrier-map
direction, the block simulation
\[
  \Theta_{CB}:\Phi_{CB}\Rightarrow\Psi_{CB}
\]
has carrier map from the \(\Psi_{CB}\)-block carrier to the
\(\Phi_{CB}\)-block carrier. The horizontal composite below is formed in the
hom-category of Mealy simulations with this displayed simulation direction.

\begin{definition}[Trace of a simulation]
\label{def:trace-simulation-feedback}
The \textbf{trace} of the simulation \(\Theta\) over \(C\) is
\[
  \operatorname{Tr}^{C}(\Theta)
  :=
  \Theta_{AB}
  \boxplus
  \coprod_{n\geq 0}
  \left(
    \Theta_{CB}
    \ast_{\mathrm h}
    \Theta_{CC}^{[n]}
    \ast_{\mathrm h}
    \Theta_{AC}
  \right),
\]
where \(\ast_{\mathrm h}\) denotes horizontal composition of simulations. The
order of \(\ast_{\mathrm h}\) follows the standing categorical convention:
\(\Theta_{AC}\) is applied first, then \(\Theta_{CC}^{[n]}\), then
\(\Theta_{CB}\).
\end{definition}

\begin{theorem}[Hom-category trace for Mealy simulations]
\label{thm:hom-category-trace-mealy-bicategory}
For internal categories \(A,B,C\), the assignment
\[
  \operatorname{Tr}^{C}_{A,B}:
  \mathbf{Mly}^{+}_{\omega}(\mathcal E)(A+C,B+C)
  \longrightarrow
  \mathbf{Mly}^{+}_{\omega}(\mathcal E)(A,B)
\]
defined on 1-cells by Definition~\ref{def:mealy-coproduct-trace} and on
2-cells by Definition~\ref{def:trace-simulation-feedback} is a functor. The
trace coherence equations are witnessed by canonical invertible simulation
2-cells.
\end{theorem}

\begin{proof}
\begin{enumerate}
  \item \textbf{Preserve identities.}
  Identity simulations restrict blockwise to identity simulations, and their
  traced homwise sum is the identity simulation on the traced Mealy morphism.

  \item \textbf{Set up vertical composition.}
  For vertical composition, let
  \[
    \Theta:\Phi\Rightarrow\Psi,
    \qquad
    \Lambda:\Psi\Rightarrow\Xi.
  \]
  By Lemma~\ref{lem:coproducts-composites-simulations-feedback}, block
  restriction commutes with vertical composition, and iterated horizontal
  composition of the \(CC\)-blocks is compatible with vertical composition by
  Lemma~\ref{lem:finite-hidden-composites-simulations}.

  \item \textbf{Compare on each hidden-length summand.}
  On the \(n\)-hidden-step summand, the traced vertical composite is
  \[
    (\Lambda\circ\Theta)_{CB}
    \ast_{\mathrm h}
    (\Lambda\circ\Theta)_{CC}^{[n]}
    \ast_{\mathrm h}
    (\Lambda\circ\Theta)_{AC}.
  \]
  By interchange for Mealy simulations, this is the vertical composite of
  \[
    \Theta_{CB}
    \ast_{\mathrm h}
    \Theta_{CC}^{[n]}
    \ast_{\mathrm h}
    \Theta_{AC}
  \]
  with
  \[
    \Lambda_{CB}
    \ast_{\mathrm h}
    \Lambda_{CC}^{[n]}
    \ast_{\mathrm h}
    \Lambda_{AC}.
  \]

  \item \textbf{Take the coproduct over hidden lengths.}
  Taking the coproduct over all \(n\geq 0\) gives
  \[
    \operatorname{Tr}^{C}(\Lambda\circ\Theta)
    =
    \operatorname{Tr}^{C}(\Lambda)\circ\operatorname{Tr}^{C}(\Theta).
  \]

  \item \textbf{Verify trace coherence.}
  The trace coherence equations are the hidden-path normal-form
  identifications from Theorem~\ref{thm:span-traced-smc}, now with a
  simulation cell decorating each edge of a hidden path. Both sides of each
  coherence equation classify the same finite decorated hidden paths.

  \item \textbf{Conclude invertibility.}
  The comparison 2-cells are induced by associativity of horizontal
  composition, associativity of pullbacks, and associativity of coproducts.
  They are invertible because the underlying hidden-path comparisons are
  canonical isomorphisms.
\end{enumerate}
\end{proof}

\section{The Int construction}
\label{sec:int-construction-feedback}

A traced symmetric monoidal category has a canonical compact closed completion,
the Int construction
\cite{JoyalStreetVerity1996TracedMonoidalCategories}. For feedback spans this
construction has a direct interface interpretation: objects become signed
interfaces, and morphisms are spans whose inputs and outputs include both
positive and negative boundary parts. This is the categorical form of the
bidirectional execution interface familiar from Geometry of Interaction
\cite{Girard1989GeometryInteractionI,AbramskyHaghverdiScott2002GeometryInteractionLCA}.

\begin{definition}[The Int category of feedback spans]
\label{def:int-feedback-spans}
The \textbf{Int category of feedback spans}
\[
  \operatorname{Int}
  \bigl(
    \mathbf{Span}^{+}_{\omega,0}(\mathcal E)
  \bigr)
\]
has objects pairs
\[
  (X^+,X^-)
\]
of objects of \(\mathcal E\). A morphism
\[
  (X^+,X^-)\longrightarrow(Y^+,Y^-)
\]
is a span in \(\mathcal E\) of type
\[
  X^+ + Y^-
  \longrightarrow
  Y^+ + X^-.
\]
Thus it is a span
\[
  X^+ + Y^-
  \xleftarrow{\ell}
  R
  \xrightarrow{r}
  Y^+ + X^-.
\]
\end{definition}

This is the standard Int convention composed with the symmetry that places the
positive target before the negative source. Thus the apparent reversal of
\(Y^-\) and \(X^-\) is only the standard contravariant negative boundary
convention, written in the chosen summand order.

Let
\[
  R:(X^+,X^-)\to(Y^+,Y^-)
\]
and
\[
  S:(Y^+,Y^-)\to(Z^+,Z^-)
\]
be Int morphisms. Thus
\[
  R:X^+ + Y^-\to Y^+ + X^-,
\]
and
\[
  S:Y^+ + Z^-\to Z^+ + Y^-.
\]
Their composite is obtained by tracing out the middle signed interface
\[
  Y^+ + Y^-.
\]
Concretely, after applying the canonical associator and symmetry spans, the
span \(R+S\) is identified with a span of type
\[
  (X^+ + Z^-)+(Y^+ + Y^-)
  \longrightarrow
  (Z^+ + X^-)+(Y^+ + Y^-).
\]
The visible source is
\[
  X^+ + Z^-,
\]
the visible target is
\[
  Z^+ + X^-,
\]
and the hidden interface is
\[
  Y^+ + Y^-.
\]
Then
\[
  S\circ_{\operatorname{Int}}R
  :=
  \operatorname{Tr}^{Y^++Y^-}(R+S),
\]
where this formula includes the preceding canonical associator and symmetry
identifications.

The two morphisms before feedback have the signed-interface shape
\[
\begin{tikzcd}[column sep=huge,row sep=large]
  X^+ + Y^-
    \ar[r,"R"]
  &
  Y^+ + X^-
  \\
  Y^+ + Z^-
    \ar[r,"S"']
  &
  Z^+ + Y^- .
\end{tikzcd}
\]
Composition feeds the \(Y^+\)- and \(Y^-\)-summands back into the middle
interface and hides them by trace:
\[
\begin{tikzcd}[column sep=large,row sep=large]
  X^+ + Z^-
    \ar[r]
  &
  (X^+ + Z^-)+(Y^+ + Y^-)
    \ar[r,"R+S"]
  &
  (Z^+ + X^-)+(Y^+ + Y^-)
    \ar[r]
  &
  Z^+ + X^- .
\end{tikzcd}
\]

\begin{theorem}[Compact closure of feedback spans]
\label{thm:int-feedback-spans}
The category
\[
  \operatorname{Int}
  \bigl(
    \mathbf{Span}^{+}_{\omega,0}(\mathcal E)
  \bigr)
\]
is compact closed.
\end{theorem}

\begin{proof}
\begin{enumerate}
  \item \textbf{Use tracedness of feedback spans.}
  By Theorem~\ref{thm:span-traced-smc},
  \[
    \bigl(\mathbf{Span}^{+}_{\omega,0}(\mathcal E),+,0\bigr)
  \]
  is traced symmetric monoidal.

  \item \textbf{Apply the Int construction.}
  The Int construction of a traced symmetric monoidal category is compact
  closed.

  \item \textbf{Specialize the description.}
  The displayed description of objects and morphisms is the standard Int
  description specialized to the coproduct tensor. Composition is the standard
  Int composition, which is precisely trace over the middle signed interface.
\end{enumerate}
\end{proof}

The same construction applies to Mealy morphisms.

\begin{definition}[The Int category of feedback Mealy morphisms]
\label{def:int-feedback-mealy}
The \textbf{Int category of feedback Mealy morphisms}
\[
  \operatorname{Int}
  \bigl(
    \mathbf{Mly}^{+}_{\omega,0}(\mathcal E)
  \bigr)
\]
has objects pairs of internal categories
\[
  (A^+,A^-).
\]
A morphism
\[
  (A^+,A^-)\longrightarrow(B^+,B^-)
\]
is a Mealy morphism
\[
  A^+ + B^-
  \rightsquigarrow
  B^+ + A^-.
\]
\end{definition}

\begin{theorem}[The Int category of feedback Mealy morphisms]
\label{thm:int-feedback-mealy}
The Int construction applied to the traced Mealy category is compact closed.
Its morphisms are signed-interface Mealy morphisms
\[
  A^+ + B^-
  \rightsquigarrow
  B^+ + A^-,
\]
and composition is given by Mealy feedback trace over the middle signed
interface.
\end{theorem}

\begin{proof}
\begin{enumerate}
  \item \textbf{Use tracedness of feedback Mealy morphisms.}
  By Theorem~\ref{thm:mly-traced-smc}, the category
  \[
    \bigl(\mathbf{Mly}^{+}_{\omega,0}(\mathcal E),+,0\bigr)
  \]
  is traced symmetric monoidal.

  \item \textbf{Apply Int.}
  Applying the Int construction gives a compact closed category.

  \item \textbf{Specialize the morphism formula.}
  The formula for morphisms is the standard Int formula specialized to
  coproduct:
  \[
    (A^+,A^-)\to(B^+,B^-)
    \quad\text{is}\quad
    A^+ + B^-\rightsquigarrow B^+ + A^-.
  \]

  \item \textbf{Identify composition.}
  Composition is trace over
  \[
    B^+ + B^-.
  \]
  The Mealy transition-output data compose by the traced Mealy construction of
  Definition~\ref{def:mealy-coproduct-trace}.
\end{enumerate}
\end{proof}

At the level of hom-categories,
Theorem~\ref{thm:hom-category-trace-mealy-bicategory} shows that simulations
are transported by the same finite hidden-path construction. Thus the Int
composition on Mealy morphisms is compatible with simulation 2-cells. A full
bicategorical Int construction would require the corresponding coherence data
for these hom-category traces.

\section{Relative execution as categorical feedback}
\label{sec:relative-execution-as-feedback}

Let
\[
  \Phi:A\rightsquigarrow B
\]
be a Mealy morphism, and let \(Y\) be a right internal \(B\)-action. Let
\[
  S_{\Phi,Y}
\]
be the relative execution-state object, and let
\[
  m_{\Phi,Y}:S_{\Phi,Y}\to S_{\Phi,Y}
\]
be the primitive execution endospan of the relative program category
\[
  \mathsf{Prog}_{\Phi,Y}.
\]
Let
\[
  i:U\hookrightarrow m_{\Phi,Y}
\]
be a sub-endospan. Put
\[
  S:=S_{\Phi,Y}.
\]
Define a span
\[
  L_U:S+S\to S+S
\]
by the block matrix
\[
  L_U=
  \begin{pmatrix}
    0_{S,S} & 1_S\\
    1_S & U
  \end{pmatrix}.
\]
The first copy of \(S\) is visible, and the second copy is hidden.

\begin{theorem}[Execution paths as feedback]
\label{thm:execution-paths-as-feedback}
There is a canonical isomorphism
\[
  \operatorname{Tr}^{S}(L_U)
  \cong
  \mathsf{Path}(U).
\]
\end{theorem}

\begin{proof}
\begin{enumerate}
  \item \textbf{Apply the trace formula.}
  By the trace formula,
  \[
    \operatorname{Tr}^{S}(L_U)
    =
    0_{S,S}
    \boxplus
    1_S\circ\mathsf{Path}(U)\circ1_S.
  \]

  \item \textbf{Remove the initial and identity factors.}
  The initial summand contributes no apex, and identity spans are units for
  span composition. Hence
  \[
    \operatorname{Tr}^{S}(L_U)
    \cong
    \mathsf{Path}(U).
  \]
\end{enumerate}
\end{proof}

Since \(m_{\Phi,Y}\) is the arrow object of an internal category over
\(S_{\Phi,Y}\), it is a monoid object in
\[
  \mathrm{EndSpan}_{\mathcal E}(S_{\Phi,Y}).
\]
The inclusion
\[
  U\hookrightarrow m_{\Phi,Y}
\]
therefore induces, by Theorem~\ref{thm:free-endospan-monoid-feedback}, a
unique monoid morphism
\[
  \operatorname{fold}_{\Phi,U,Y}:
  \mathsf{Path}(U)\longrightarrow m_{\Phi,Y}.
\]

\begin{corollary}[Folded feedback]
\label{cor:folded-feedback}
The folded execution map factors canonically as
\[
  \operatorname{Tr}^{S}(L_U)
  \cong
  \mathsf{Path}(U)
  \xrightarrow{\operatorname{fold}_{\Phi,U,Y}}
  m_{\Phi,Y}.
\]
\end{corollary}

\section{Predicate transformers for feedback spans}
\label{sec:predicate-transformers-feedback}

We now add the logical structure needed for modalities. The categorical
implementation uses subobjects as predicates, regular images as existential
quantification, and right adjoints to pullback as universal quantification
\cite{Jacobs1999CategoricalLogic}. In regular and relational settings this is
the standard categorical form of relation and program semantics
\cite{FongSpivakRegularRelational,HarelKozenTiurynDynamicLogic}.

\begin{definition}[Modal feedback base]
\label{def:modal-feedback-base}
A \textbf{modal feedback base} is a feedback span base \(\mathcal E\) whose
countable coproducts are subobject-extensive and whose subobject fibration has
the regular, first-order Heyting, and \(\omega\)-iterative structure used for
dynamic span logic. Thus existential images of subobjects exist and are stable
under pullback; pullback on subobject fibres has both adjoints
\[
  \exists_f\dashv f^\ast\dashv\forall_f;
\]
each relevant subobject fibre has the countable joins and countable meets used
below; and the pullback functors \(f^\ast\) preserve those countable joins and
meets.

In particular, for every relevant countable coproduct
\[
  A\cong\coprod_{n\geq 0}A_n
\]
the canonical comparison
\[
  \operatorname{Sub}(A)
  \longrightarrow
  \prod_{n\geq 0}\operatorname{Sub}(A_n)
\]
is an isomorphism.
\end{definition}

Assume throughout the remainder of the logical sections that \(\mathcal E\) is
a modal feedback base.

Let
\[
  R:X\to Y
\]
be a span
\[
  X\xleftarrow{\ell}R\xrightarrow{r}Y.
\]
For
\[
  \varphi\in\operatorname{Sub}(Y),
\]
define
\[
  \langle R\rangle\varphi
  :=
  \exists_{\ell}(r^\ast\varphi),
\]
and
\[
  [R]\varphi
  :=
  \forall_{\ell}(r^\ast\varphi).
\]

The definition is encoded by the pullback square
\[
\begin{tikzcd}[column sep=large,row sep=large]
  R\times_Y \varphi
    \ar[r]
    \ar[d]
  &
  \varphi
    \ar[d,hook]
  \\
  R
    \ar[r,"r"']
  &
  Y
\end{tikzcd}
\qquad
\text{followed by}
\qquad
R\times_Y\varphi\to R\xrightarrow{\ell}X .
\]

\begin{lemma}[Modal calculus for spans]
\label{lem:modal-calculus-feedback}
Let \(R,S:X\to Y\) be parallel spans, and let
\[
  R:X\to Y,
  \qquad
  T:Y\to Z
\]
be composable spans. Then
\[
  \langle R\boxplus S\rangle\varphi
  =
  \langle R\rangle\varphi
  \vee
  \langle S\rangle\varphi,
\]
\[
  [R\boxplus S]\varphi
  =
  [R]\varphi
  \wedge
  [S]\varphi,
\]
and
\[
  \langle T\circ R\rangle\varphi
  =
  \langle R\rangle\langle T\rangle\varphi,
\]
\[
  [T\circ R]\varphi
  =
  [R][T]\varphi.
\]
\end{lemma}

\begin{proof}
\begin{enumerate}
  \item \textbf{Compute modalities of homwise coproducts.}
  For homwise coproducts, the left leg of \(R\boxplus S\) is the coproduct of
  the left legs of \(R\) and \(S\). Existential quantification along a
  coproduct of maps gives the join of the existential images, while universal
  quantification along a coproduct of maps gives the meet of the universal
  images. The subobject-extensive condition identifies the relevant predicates
  on coproduct apexes with summandwise predicates.

  \item \textbf{Compute modalities of composites.}
  For composition, the apex of \(T\circ R\) is the pullback
  \[
    R\times_Y T.
  \]
  The Beck--Chevalley isomorphism for this pullback square,
  \[
  \begin{tikzcd}[column sep=large,row sep=large]
    R\times_Y T
      \ar[r]
      \ar[d]
    &
    T
      \ar[d]
    \\
    R
      \ar[r]
    &
    Y ,
  \end{tikzcd}
  \]
  together with
  \[
    \exists\dashv(-)^\ast\dashv\forall,
  \]
  gives the composite predicate-transformer identities.
\end{enumerate}
\end{proof}

\begin{lemma}[Tests]
\label{lem:test-modalities-feedback}
Let
\[
  \theta:\Theta\hookrightarrow S
\]
be a subobject, and let
\[
  T_\theta:=
  \left(
    S\xleftarrow{\theta}\Theta\xrightarrow{\theta}S
  \right)
\]
be its test span. Then
\[
  \langle T_\theta\rangle\varphi
  =
  \theta\wedge\varphi,
\]
and
\[
  [T_\theta]\varphi
  =
  \theta\Rightarrow\varphi.
\]
\end{lemma}

\begin{proof}
\begin{enumerate}
  \item \textbf{Compute the diamond.}
  For diamonds,
  \[
    \langle T_\theta\rangle\varphi
    =
    \exists_\theta(\theta^\ast\varphi).
  \]
  Since \(\theta\) is monic, this is the image of the pullback intersection
  of \(\theta\) and \(\varphi\), hence
  \[
    \langle T_\theta\rangle\varphi
    =
    \theta\wedge\varphi.
  \]

  \item \textbf{Compute the box by adjunction.}
  For boxes,
  \[
    [T_\theta]\varphi
    =
    \forall_\theta(\theta^\ast\varphi).
  \]
  For any predicate \(\psi\hookrightarrow S\),
  \[
    \psi\leq\forall_\theta(\theta^\ast\varphi)
  \]
  if and only if
  \[
    \theta^\ast\psi\leq\theta^\ast\varphi.
  \]
  Since \(\theta\) is monic, this is equivalent to
  \[
    \psi\wedge\theta\leq\varphi.
  \]
  By the Heyting adjunction, this is equivalent to
  \[
    \psi\leq\theta\Rightarrow\varphi.
  \]
  Therefore
  \[
    [T_\theta]\varphi
    =
    \theta\Rightarrow\varphi.
  \]
\end{enumerate}
\end{proof}

\begin{lemma}[Modalities of countable homwise coproducts]
\label{lem:countable-coproduct-modalities-feedback}
Let
\[
  (R_n)_{n\geq 0}:X\to Y
\]
be a countable family of parallel spans. Then
\[
  \left\langle
    \boxplus_{n\geq 0}R_n
  \right\rangle\varphi
  =
  \bigvee_{n\geq 0}
  \langle R_n\rangle\varphi,
\]
and
\[
  \left[
    \boxplus_{n\geq 0}R_n
  \right]\varphi
  =
  \bigwedge_{n\geq 0}
  [R_n]\varphi.
\]
\end{lemma}

\begin{proof}
\begin{enumerate}
  \item \textbf{Use subobject-extensivity.}
  By subobject-extensivity, a subobject of the coproduct apex
  \[
    \coprod_{n\geq 0}R_n
  \]
  is determined by its restrictions to the summands.

  \item \textbf{Compute the diamond.}
  Existential image along the coproduct source leg is the join of the
  summandwise existential images. Hence
  \[
    \left\langle
      \boxplus_{n\geq 0}R_n
    \right\rangle\varphi
    =
    \bigvee_{n\geq 0}
    \langle R_n\rangle\varphi.
  \]

  \item \textbf{Compute the box by testing against a predicate.}
  Let
  \[
    L:\coprod_{n\geq 0}R_n\to X
  \]
  be the coproduct source leg, and let
  \[
    R:\coprod_{n\geq 0}R_n\to Y
  \]
  be the coproduct target leg. For any
  \[
    \psi\hookrightarrow X,
  \]
  one has
  \[
    \psi\leq
    \left[
      \boxplus_{n\geq 0}R_n
    \right]\varphi
  \]
  if and only if
  \[
    L^\ast\psi\leq R^\ast\varphi.
  \]

  \item \textbf{Restrict to summands.}
  By subobject-extensivity, this is equivalent to the family of inequalities
  \[
    \ell_n^\ast\psi\leq r_n^\ast\varphi
    \qquad(n\geq 0),
  \]
  where \(\ell_n,r_n\) are the legs of \(R_n\).

  \item \textbf{Use the adjunctions.}
  By the adjunction
  \[
    \ell_n^\ast\dashv\forall_{\ell_n},
  \]
  these inequalities are equivalent to
  \[
    \psi\leq [R_n]\varphi
    \qquad(n\geq 0).
  \]
  Hence
  \[
    \psi\leq
    \bigwedge_{n\geq 0}[R_n]\varphi.
  \]
  Therefore
  \[
    \left[
      \boxplus_{n\geq 0}R_n
    \right]\varphi
    =
    \bigwedge_{n\geq 0}
    [R_n]\varphi.
  \]
\end{enumerate}
\end{proof}

\section{Fixed-point semantics of finite paths}
\label{sec:fixed-point-semantics-path}

\begin{theorem}[Finite-path fixed-point semantics]
\label{thm:finite-path-fixed-point-semantics}
For every endospan \(U:S\to S\) and every
\[
  \varphi\in\operatorname{Sub}(S),
\]
one has
\[
  \langle\mathsf{Path}(U)\rangle\varphi
  =
  \mu X.\bigl(\varphi\vee\langle U\rangle X\bigr),
\]
and
\[
  [\mathsf{Path}(U)]\varphi
  =
  \nu X.\bigl(\varphi\wedge[U]X\bigr).
\]
\end{theorem}

\begin{proof}
\begin{enumerate}
  \item \textbf{Expand the diamond of a path object.}
  By Lemma~\ref{lem:countable-coproduct-modalities-feedback}, for diamonds,
  \[
    \langle\mathsf{Path}(U)\rangle\varphi
    =
    \left\langle\coprod_{n\geq 0}U^{[n]}\right\rangle\varphi
    =
    \bigvee_{n\geq 0}\langle U^{[n]}\rangle\varphi.
  \]
  Let
  \[
    P:=
    \bigvee_{n\geq 0}\langle U^{[n]}\rangle\varphi.
  \]

  \item \textbf{Show that \(P\) is a fixed point.}
  Since \(U^{[0]}=1_S\), we have \(\varphi\leq P\). Since
  \[
    U^{[n+1]}=U^{[n]}\circ U,
  \]
  we have
  \[
    \langle U^{[n+1]}\rangle\varphi
    =
    \langle U\rangle\langle U^{[n]}\rangle\varphi.
  \]
  The diamond transformer \(\langle U\rangle\) preserves the joins used here,
  so
  \[
    \langle U\rangle P
    =
    \bigvee_{n\geq 0}\langle U^{[n+1]}\rangle\varphi
    \leq P.
  \]
  Thus
  \[
    \varphi\vee\langle U\rangle P\leq P.
  \]
  Conversely,
  \[
    P
    =
    \varphi
    \vee
    \bigvee_{n\geq 0}\langle U^{[n+1]}\rangle\varphi
    =
    \varphi\vee\langle U\rangle P.
  \]

  \item \textbf{Show leastness.}
  If \(Q\) is any pre-fixed point of
  \[
    F(X):=\varphi\vee\langle U\rangle X,
  \]
  then
  \[
    \varphi\vee\langle U\rangle Q\leq Q.
  \]
  Hence \(\varphi\leq Q\) and \(\langle U\rangle Q\leq Q\). By induction,
  \[
    \langle U^{[n]}\rangle\varphi\leq Q
  \]
  for every \(n\). Taking the join over \(n\) gives \(P\leq Q\). Therefore
  \(P=\mu F\).

  \item \textbf{Expand the box of a path object.}
  By Lemma~\ref{lem:countable-coproduct-modalities-feedback}, for boxes,
  \[
    [\mathsf{Path}(U)]\varphi
    =
    \left[\coprod_{n\geq 0}U^{[n]}\right]\varphi
    =
    \bigwedge_{n\geq 0}[U^{[n]}]\varphi.
  \]
  Let
  \[
    Q:=
    \bigwedge_{n\geq 0}[U^{[n]}]\varphi.
  \]

  \item \textbf{Show that \(Q\) is a fixed point.}
  Since \(U^{[0]}=1_S\), we have \(Q\leq\varphi\). Since
  \[
    [U^{[n+1]}]\varphi
    =
    [U][U^{[n]}]\varphi,
  \]
  and \([U]\) preserves the meets used here,
  \[
    [U]Q
    =
    \bigwedge_{n\geq 0}[U^{[n+1]}]\varphi.
  \]
  Hence
  \[
    Q\leq[U]Q,
  \]
  and so
  \[
    Q\leq\varphi\wedge[U]Q.
  \]
  Conversely,
  \[
    \varphi\wedge[U]Q
    =
    [U^{[0]}]\varphi
    \wedge
    \bigwedge_{n\geq 0}[U^{[n+1]}]\varphi
    =
    Q.
  \]

  \item \textbf{Show greatestness.}
  If \(P\) is any post-fixed point of
  \[
    G(X):=\varphi\wedge[U]X,
  \]
  then
  \[
    P\leq\varphi\wedge[U]P.
  \]
  Hence \(P\leq\varphi\) and \(P\leq[U]P\). By induction,
  \[
    P\leq[U^{[n]}]\varphi
  \]
  for every \(n\). Taking the meet over \(n\) gives \(P\leq Q\). Therefore
  \(Q=\nu G\).
\end{enumerate}
\end{proof}

\begin{theorem}[Modal semantics of the dagger]
\label{thm:modal-semantics-dagger-feedback}
For every span
\[
  f:X\to Y+X
\]
and every
\[
  \varphi\in\operatorname{Sub}(Y),
\]
one has
\[
  \langle f^\dagger\rangle\varphi
  =
  \mu Z.
  \bigl(
    \langle f_Y\rangle\varphi
    \vee
    \langle f_X\rangle Z
  \bigr),
\]
and
\[
  [f^\dagger]\varphi
  =
  \nu Z.
  \bigl(
    [f_Y]\varphi
    \wedge
    [f_X]Z
  \bigr).
\]
\end{theorem}

\begin{proof}
\begin{enumerate}
  \item \textbf{Compute the diamond.}
  Using
  \[
    f^\dagger=f_Y\circ\mathsf{Path}(f_X),
  \]
  the diamond calculation is
  \[
  \begin{aligned}
    \langle f^\dagger\rangle\varphi
    &=
    \langle\mathsf{Path}(f_X)\rangle
    \langle f_Y\rangle\varphi
    \\
    &=
    \mu Z.
    \bigl(
      \langle f_Y\rangle\varphi
      \vee
      \langle f_X\rangle Z
    \bigr),
  \end{aligned}
  \]
  by Theorem~\ref{thm:finite-path-fixed-point-semantics}.

  \item \textbf{Compute the box.}
  The box calculation is dual:
  \[
  \begin{aligned}
    [f^\dagger]\varphi
    &=
    [\mathsf{Path}(f_X)]
    [f_Y]\varphi
    \\
    &=
    \nu Z.
    \bigl(
      [f_Y]\varphi
      \wedge
      [f_X]Z
    \bigr).
  \end{aligned}
  \]
\end{enumerate}
\end{proof}

\section{Dynamic logic of feedback trace}
\label{sec:dynamic-logic-categorical-feedback}

Let
\[
  R:X+H\to Y+H
\]
have block decomposition
\[
  R=
  \begin{pmatrix}
    R_{XY} & R_{XH}\\
    R_{HY} & R_{HH}
  \end{pmatrix}.
\]

\begin{theorem}[Diamond law for feedback trace]
\label{thm:diamond-law-feedback-trace}
For every
\[
  \varphi\in\operatorname{Sub}(Y),
\]
one has
\[
  \left\langle
    \operatorname{Tr}^{H}(R)
  \right\rangle\varphi
  =
  \langle R_{XY}\rangle\varphi
  \vee
  \langle R_{XH}\rangle
  \mu Z.
  \bigl(
    \langle R_{HY}\rangle\varphi
    \vee
    \langle R_{HH}\rangle Z
  \bigr).
\]
\end{theorem}

\begin{proof}
\begin{enumerate}
  \item \textbf{Expand the trace formula.}
  Using the trace formula and Lemma~\ref{lem:modal-calculus-feedback},
  \[
  \begin{aligned}
    \left\langle
      \operatorname{Tr}^{H}(R)
    \right\rangle\varphi
    &=
    \left\langle
      R_{XY}
      \boxplus
      R_{HY}\circ\mathsf{Path}(R_{HH})\circ R_{XH}
    \right\rangle\varphi
    \\
    &=
    \langle R_{XY}\rangle\varphi
    \vee
    \langle R_{XH}\rangle
    \langle\mathsf{Path}(R_{HH})\rangle
    \langle R_{HY}\rangle\varphi.
  \end{aligned}
  \]

  \item \textbf{Apply finite-path fixed-point semantics.}
  By Theorem~\ref{thm:finite-path-fixed-point-semantics},
  \[
    \langle\mathsf{Path}(R_{HH})\rangle
    \langle R_{HY}\rangle\varphi
    =
    \mu Z.
    \bigl(
      \langle R_{HY}\rangle\varphi
      \vee
      \langle R_{HH}\rangle Z
    \bigr).
  \]

  \item \textbf{Substitute the fixed point.}
  Substitution gives the displayed formula.
\end{enumerate}
\end{proof}

\begin{theorem}[Box law for feedback trace]
\label{thm:box-law-feedback-trace}
For every
\[
  \varphi\in\operatorname{Sub}(Y),
\]
one has
\[
  \left[
    \operatorname{Tr}^{H}(R)
  \right]\varphi
  =
  [R_{XY}]\varphi
  \wedge
  [R_{XH}]
  \nu Z.
  \bigl(
    [R_{HY}]\varphi
    \wedge
    [R_{HH}]Z
  \bigr).
\]
\end{theorem}

\begin{proof}
\begin{enumerate}
  \item \textbf{Expand the trace formula.}
  The box modality sends homwise coproducts of spans to meets and categorical
  composition to composition of box predicate transformers. Therefore
  \[
  \begin{aligned}
    \left[
      \operatorname{Tr}^{H}(R)
    \right]\varphi
    &=
    \left[
      R_{XY}
      \boxplus
      R_{HY}\circ\mathsf{Path}(R_{HH})\circ R_{XH}
    \right]\varphi
    \\
    &=
    [R_{XY}]\varphi
    \wedge
    [R_{XH}]
    [\mathsf{Path}(R_{HH})]
    [R_{HY}]\varphi.
  \end{aligned}
  \]

  \item \textbf{Apply finite-path fixed-point semantics.}
  By Theorem~\ref{thm:finite-path-fixed-point-semantics},
  \[
    [\mathsf{Path}(R_{HH})]
    [R_{HY}]\varphi
    =
    \nu Z.
    \bigl(
      [R_{HY}]\varphi
      \wedge
      [R_{HH}]Z
    \bigr).
  \]

  \item \textbf{Substitute the fixed point.}
  Substitution gives the displayed formula.
\end{enumerate}
\end{proof}

\section{Typed dynamic program syntax}
\label{sec:dynamic-program-syntax-feedback}

Program terms are typed. A judgement
\[
  p:S\leadsto T
\]
means that \(p\) denotes a span from \(S\) to \(T\). The symbol
\(\leadsto\) is used here for program typing, to distinguish program terms
from Mealy morphisms. The syntax follows the standard dynamic-logic program
operations of sequencing, choice, iteration, and tests
\cite{HarelKozenTiurynDynamicLogic}.

Primitive programs are supplied by a chosen family of primitive spans. If
\[
  r:S\leadsto T
\]
is primitive, then
\[
  \llbracket r\rrbracket:S\to T
\]
is its corresponding span.

The compound term formers are:
\[
  \frac{p:S\leadsto T \qquad q:T\leadsto U}
       {p;q:S\leadsto U},
\]
\[
  \frac{p:S\leadsto T \qquad q:S\leadsto T}
       {p\cup q:S\leadsto T},
\]
\[
  \frac{p:S\leadsto S}
       {p^{\star_{\mathrm{path}}}:S\leadsto S},
\]
and
\[
  \frac{\theta:\Theta\hookrightarrow S}
       {\theta?:S\leadsto S}.
\]

The interpretation is defined by
\[
  \llbracket p;q\rrbracket
  =
  \llbracket q\rrbracket\circ\llbracket p\rrbracket,
\]
\[
  \llbracket p\cup q\rrbracket
  =
  \llbracket p\rrbracket\boxplus\llbracket q\rrbracket,
\]
\[
  \llbracket p^{\star_{\mathrm{path}}}\rrbracket
  =
  \mathsf{Path}(\llbracket p\rrbracket).
\]
For a predicate
\[
  \theta:\Theta\hookrightarrow S,
\]
define its test span
\[
  T_{\theta}:=
  \left(
    S\xleftarrow{\theta}\Theta\xrightarrow{\theta}S
  \right),
\]
and set
\[
  \llbracket\theta?\rrbracket:=T_{\theta}.
\]

The modal semantics of program terms is
\[
  \langle p\rangle\varphi
  :=
  \langle\llbracket p\rrbracket\rangle\varphi,
  \qquad
  [p]\varphi
  :=
  [\llbracket p\rrbracket]\varphi.
\]
Sequential composition and iteration are therefore syntactic program-forming
operations whose categorical interpretations are span composition and the
finite-path construction.

\section{Guarded trace and guarded while programs}
\label{sec:guarded-trace-and-while-feedback}

Let
\[
  \theta:\Theta\hookrightarrow H
\]
be a predicate on \(H\), and let
\[
  T_\theta:H\to H
\]
be its test span. Abstract guarded traces are studied in
\cite{GoncharovSchroder2018GuardedTracedCategories}. The construction below
is a predicate-guarded specialization inside the total traced span category
constructed above.

\begin{definition}[Source-guarded endospan]
\label{def:predicate-guarded-endospan}
An endospan
\[
  U:H\to H
\]
is \textbf{source \(\theta\)-guarded} if its left leg factors through
\(\theta\). Equivalently, because \(\theta\) is monic, the canonical
comparison
\[
  U\circ T_\theta\longrightarrow U
\]
is an isomorphism of endospans.
\end{definition}

The guard condition is the factorization of the source leg:
\[
\begin{tikzcd}[column sep=large,row sep=large]
  U
    \ar[rr,"\bar\ell"]
    \ar[dr,"\ell"']
  &&
  \Theta
    \ar[dl,hook,"\theta"]
  \\
  &
  H .
\end{tikzcd}
\]

\begin{definition}[Predicate-guarded feedback situation]
\label{def:predicate-guarded-trace}
Let
\[
  R:X+H\to Y+H
\]
be a span whose hidden block
\[
  R_{HH}:H\to H
\]
is source \(\theta\)-guarded. We call \(R\) a
\textbf{\(\theta\)-guarded feedback situation}. Its feedback trace is the
ordinary coproduct trace
\[
  \operatorname{Tr}^{H}_{X,Y}(R)
  =
  R_{XY}
  \boxplus
  R_{HY}\circ\mathsf{Path}(R_{HH})\circ R_{XH}.
\]
\end{definition}

\begin{proposition}[Guarded hidden-path normal form]
\label{prop:guarded-hidden-path-normal-form}
If \(R\) is a \(\theta\)-guarded feedback situation, then each positive
occurrence of \(R_{HH}\) in a hidden path may be canonically rewritten by
inserting \(T_\theta\) on its source side:
\[
  R_{HH}\cong R_{HH}\circ T_\theta .
\]
Equivalently, in categorical notation each occurrence of \(R_{HH}\) may be
replaced by \(R_{HH}\circ T_\theta\), so operationally the guard is tested
immediately before that hidden step is taken. Hence the trace is the homwise
coproduct of finite hidden paths whose positive hidden-loop steps are guarded
at their source by \(\theta\).
\end{proposition}

\begin{proof}
\begin{enumerate}
  \item \textbf{Use source-guardedness.}
  Since \(R_{HH}\) is source \(\theta\)-guarded, the canonical comparison
  \[
    R_{HH}\circ T_\theta\longrightarrow R_{HH}
  \]
  is an isomorphism.

  \item \textbf{Rewrite each hidden loop step.}
  Therefore each occurrence of \(R_{HH}\) in an iterated hidden path may be
  written as an occurrence of \(R_{HH}\circ T_\theta\). The \(n\)-hidden-step
  summand of the trace is
  \[
    R_{HY}\circ R_{HH}^{[n]}\circ R_{XH}.
  \]

  \item \textbf{Insert the source guards.}
  For \(n>0\), this composite contains \(n\) occurrences of \(R_{HH}\), and
  each such occurrence may be preceded operationally by \(T_\theta\), i.e.
  replaced categorically by \(R_{HH}\circ T_\theta\).

  \item \textbf{Take the coproduct over finite lengths.}
  Taking the coproduct over \(n\geq 0\) gives the claimed guarded hidden-path
  normal form.
\end{enumerate}
\end{proof}

Let \(u:S\leadsto S\) be a program with interpretation
\[
  U:=\llbracket u\rrbracket:S\to S.
\]
Let
\[
  \theta:\Theta\hookrightarrow S,
  \qquad
  \epsilon:E\hookrightarrow S
\]
be loop and exit predicates. Define a block span
\[
  W_{\theta,\epsilon}(U):S+S\to S+S
\]
by
\[
  W_{\theta,\epsilon}(U)
  =
  \begin{pmatrix}
    0_{S,S} & 1_S\\
    T_\epsilon & U\circ T_\theta
  \end{pmatrix}.
\]
The first copy of \(S\) is visible, and the second copy is hidden. The hidden
loop block
\[
  U\circ T_\theta
\]
is source \(\theta\)-guarded.

\begin{theorem}[Guarded while as guarded feedback]
\label{thm:guarded-while-as-guarded-trace}
There is a canonical isomorphism
\[
  \operatorname{Tr}^{S}
  \bigl(
    W_{\theta,\epsilon}(U)
  \bigr)
  \cong
  T_\epsilon\circ\mathsf{Path}(U\circ T_\theta).
\]
Equivalently,
\[
  \operatorname{Tr}^{S}
  \bigl(
    W_{\theta,\epsilon}(U)
  \bigr)
  \cong
  \llbracket
    (\theta?;u)^{\star_{\mathrm{path}}};\epsilon?
  \rrbracket .
\]
\end{theorem}

\begin{proof}
\begin{enumerate}
  \item \textbf{Apply the feedback trace formula.}
  By the feedback trace formula,
  \[
    \operatorname{Tr}^{S}
    \bigl(
      W_{\theta,\epsilon}(U)
    \bigr)
    =
    0_{S,S}
    \boxplus
    T_\epsilon
    \circ
    \mathsf{Path}(U\circ T_\theta)
    \circ
    1_S.
  \]

  \item \textbf{Remove the initial and identity factors.}
  Removing the initial summand and the identity span gives
  \[
    T_\epsilon\circ\mathsf{Path}(U\circ T_\theta).
  \]

  \item \textbf{Identify the program interpretation.}
  For the syntactic expression,
  \[
    \llbracket\theta?;u\rrbracket
    =
    \llbracket u\rrbracket\circ\llbracket\theta?\rrbracket
    =
    U\circ T_\theta.
  \]
  Therefore
  \[
    \llbracket
      (\theta?;u)^{\star_{\mathrm{path}}};\epsilon?
    \rrbracket
    =
    T_\epsilon\circ\mathsf{Path}(U\circ T_\theta).
  \]
  This is represented by
  \[
  \begin{tikzcd}[column sep=large,row sep=large]
    S
      \ar[r,"T_\theta"]
    &
    S
      \ar[r,"U"]
    &
    S
      \ar[r,"T_\theta"]
    &
    S
      \ar[r,"U"]
    &
    S
      \ar[r,"\cdots"]
    &
    S
      \ar[r,"T_\epsilon"]
    &
    S .
  \end{tikzcd}
  \]
\end{enumerate}
\end{proof}

\begin{theorem}[Diamond law for guarded while]
\label{thm:guarded-while-diamond-feedback}
For every
\[
  \varphi\in\operatorname{Sub}(S),
\]
one has
\[
  \left\langle
    (\theta?;u)^{\star_{\mathrm{path}}};\epsilon?
  \right\rangle\varphi
  =
  \mu X.
  \bigl(
    \epsilon\wedge\varphi
    \vee
    \theta\wedge\langle u\rangle X
  \bigr).
\]
\end{theorem}

\begin{proof}
\begin{enumerate}
  \item \textbf{Use the guarded while interpretation.}
  Let \(U=\llbracket u\rrbracket\). By
  Theorem~\ref{thm:guarded-while-as-guarded-trace},
  \[
    \llbracket
      (\theta?;u)^{\star_{\mathrm{path}}};\epsilon?
    \rrbracket
    =
    T_\epsilon\circ\mathsf{Path}(U\circ T_\theta).
  \]
  Therefore
  \[
    \left\langle
      (\theta?;u)^{\star_{\mathrm{path}}};\epsilon?
    \right\rangle\varphi
    =
    \langle\mathsf{Path}(U\circ T_\theta)\rangle
    \langle T_\epsilon\rangle\varphi.
  \]

  \item \textbf{Evaluate the tests.}
  By Lemma~\ref{lem:test-modalities-feedback},
  \[
    \langle T_\epsilon\rangle\varphi
    =
    \epsilon\wedge\varphi,
  \]
  and
  \[
    \langle U\circ T_\theta\rangle X
    =
    \langle T_\theta\rangle\langle U\rangle X
    =
    \theta\wedge\langle u\rangle X.
  \]

  \item \textbf{Apply finite-path fixed-point semantics.}
  Applying Theorem~\ref{thm:finite-path-fixed-point-semantics} to
  \(U\circ T_\theta\) gives the displayed formula.
\end{enumerate}
\end{proof}

\begin{theorem}[Box law for guarded while]
\label{thm:guarded-while-box-feedback}
For every
\[
  \varphi\in\operatorname{Sub}(S),
\]
one has
\[
  \left[
    (\theta?;u)^{\star_{\mathrm{path}}};\epsilon?
  \right]\varphi
  =
  \nu X.
  \bigl(
    (\epsilon\Rightarrow\varphi)
    \wedge
    (\theta\Rightarrow[u]X)
  \bigr).
\]
\end{theorem}

\begin{proof}
\begin{enumerate}
  \item \textbf{Use the guarded while interpretation.}
  Let \(U=\llbracket u\rrbracket\). By
  Theorem~\ref{thm:guarded-while-as-guarded-trace},
  \[
    \llbracket
      (\theta?;u)^{\star_{\mathrm{path}}};\epsilon?
    \rrbracket
    =
    T_\epsilon\circ\mathsf{Path}(U\circ T_\theta).
  \]
  Therefore
  \[
    \left[
      (\theta?;u)^{\star_{\mathrm{path}}};\epsilon?
    \right]\varphi
    =
    [\mathsf{Path}(U\circ T_\theta)]
    [T_\epsilon]\varphi.
  \]

  \item \textbf{Evaluate the tests.}
  By Lemma~\ref{lem:test-modalities-feedback},
  \[
    [T_\epsilon]\varphi
    =
    \epsilon\Rightarrow\varphi,
  \]
  and
  \[
    [U\circ T_\theta]X
    =
    [T_\theta][U]X
    =
    \theta\Rightarrow[u]X.
  \]

  \item \textbf{Apply finite-path fixed-point semantics.}
  Applying Theorem~\ref{thm:finite-path-fixed-point-semantics} to
  \(U\circ T_\theta\) gives the displayed formula.
\end{enumerate}
\end{proof}

\section{Stability under downstream change and simulations}
\label{sec:feedback-stability}

The feedback construction is stable under the source-cartesian change of
execution spans induced by downstream action morphisms and by Mealy
simulations.

\begin{definition}[Source-cartesian comparison of execution endospans]
\label{def:source-cartesian-comparison-execution-endospans}
Let
\[
  U=
  \left(
    S\xleftarrow{\ell_U}U\xrightarrow{r_U}S
  \right),
  \qquad
  V=
  \left(
    S'\xleftarrow{\ell_V}V\xrightarrow{r_V}S'
  \right)
\]
be endospans. A \textbf{source-cartesian comparison}
\[
  (k,h):U\longrightarrow V
\]
consists of maps
\[
  h:S\to S',
  \qquad
  k:U\to V
\]
such that
\[
  \ell_V k=h\ell_U,
  \qquad
  r_V k=hr_U,
\]
and the square
\[
\begin{tikzcd}[column sep=large,row sep=large]
  U
    \ar[r,"k"]
    \ar[d,"\ell_U"']
  &
  V
    \ar[d,"\ell_V"]
  \\
  S
    \ar[r,"h"']
  &
  S'
\end{tikzcd}
\]
is a pullback.
\end{definition}

Equivalently, \(U\) is obtained by pulling back the source leg of \(V\) along
\(h\), while the target leg is transported by the equation \(r_Vk=hr_U\).
This is the form of comparison used by execution functors induced by
downstream changes and by Mealy simulations.

For a Mealy simulation
\[
  \Theta:\Phi\Rightarrow\Psi,
\]
the comparison map is taken in the direction of the induced program functor
\[
  h_\Theta:S_{\Psi,Y}\to S_{\Phi,Y}.
\]
Thus the primitive span for \(\Psi\) is compared source-cartesianly with the
primitive span for \(\Phi\) along \(h_\Theta\). This fixes the direction of
the simulation-induced stability statement.

\begin{lemma}[Primitive modal stability for source-cartesian comparisons]
\label{lem:primitive-source-cartesian-modal-stability}
Let
\[
  (k,h):U\to V
\]
be a source-cartesian comparison of endospans as in
Definition~\ref{def:source-cartesian-comparison-execution-endospans}. For
every
\[
  \psi\in\operatorname{Sub}(S'),
\]
one has
\[
  h^\ast\langle V\rangle\psi
  =
  \langle U\rangle h^\ast\psi,
\]
and
\[
  h^\ast[V]\psi
  =
  [U]h^\ast\psi.
\]
\end{lemma}

\begin{proof}
\begin{enumerate}
  \item \textbf{Use the source pullback square.}
  The defining pullback square
  \[
  \begin{tikzcd}[column sep=large,row sep=large]
    U
      \ar[r,"k"]
      \ar[d,"\ell_U"']
    &
    V
      \ar[d,"\ell_V"]
    \\
    S
      \ar[r,"h"']
    &
    S'
  \end{tikzcd}
  \]
  gives Beck--Chevalley isomorphisms for existential and universal
  quantification along the source legs.

  \item \textbf{Use compatibility of the target legs.}
  The equality
  \[
    r_Vk=hr_U
  \]
  identifies
  \[
    k^\ast r_V^\ast\psi
    =
    r_U^\ast h^\ast\psi.
  \]

  \item \textbf{Compute the diamond.}
  Therefore
  \[
  \begin{aligned}
    h^\ast\langle V\rangle\psi
    &=
    h^\ast\exists_{\ell_V}(r_V^\ast\psi)
    \\
    &=
    \exists_{\ell_U}k^\ast(r_V^\ast\psi)
    \\
    &=
    \exists_{\ell_U}(r_U^\ast h^\ast\psi)
    \\
    &=
    \langle U\rangle h^\ast\psi.
  \end{aligned}
  \]

  \item \textbf{Compute the box.}
  The box calculation is the dual Beck--Chevalley calculation:
  \[
  \begin{aligned}
    h^\ast[V]\psi
    &=
    h^\ast\forall_{\ell_V}(r_V^\ast\psi)
    \\
    &=
    \forall_{\ell_U}k^\ast(r_V^\ast\psi)
    \\
    &=
    \forall_{\ell_U}(r_U^\ast h^\ast\psi)
    \\
    &=
    [U]h^\ast\psi.
  \end{aligned}
  \]
\end{enumerate}
\end{proof}

\begin{lemma}[Iterated source-cartesian comparisons]
\label{lem:iterated-source-cartesian-comparisons}
Let
\[
  (k,h):U\to V
\]
be a source-cartesian comparison of endospans. For every \(n\geq 0\), there is
an induced source-cartesian comparison
\[
  (k^{[n]},h):U^{[n]}\longrightarrow V^{[n]}.
\]
\end{lemma}

\begin{proof}
\begin{enumerate}
  \item \textbf{Handle length zero.}
  For \(n=0\), both endospans are identity spans. The comparison is the
  identity comparison over \(h\), and the source square is the canonical
  pullback square for identities.

  \item \textbf{Handle length one.}
  For \(n=1\), this is the given comparison
  \[
    (k,h):U\to V.
  \]

  \item \textbf{Use induction and the pullback cube for span composition.}
  Suppose the comparison exists for \(n\). The apex of
  \[
    U^{[n+1]}=U^{[n]}\circ U
  \]
  is
  \[
    U\times_S U^{[n]}.
  \]
  The comparison apex map is
  \[
    k^{[n+1]}:=k\times_h k^{[n]}:
    U\times_S U^{[n]}
    \longrightarrow
    V\times_{S'}V^{[n]}.
  \]
  The source leg of \(U^{[n+1]}\) is the source leg of the rightmost
  operational \(U\)-factor. Since the square defining \((k,h):U\to V\) is a
  pullback, and since the second factor is matched by the target-compatibility
  equation
  \[
    r_Vk=hr_U,
  \]
  the standard pullback-pasting/cube argument identifies
  \[
    U\times_S U^{[n]}
    \cong
    S\times_{S'}(V\times_{S'}V^{[n]}).
  \]
  Thus the source square for
  \[
    (k^{[n+1]},h):U^{[n+1]}\to V^{[n+1]}
  \]
  is again a pullback.

  \item \textbf{Preserve target compatibility.}
  The target-leg equality for the composite follows from the target-leg
  equalities for the two factors and associativity of span composition.

  \item \textbf{Conclude source-cartesianness.}
  Pullbacks are closed under pasting, so the source square for the composite
  is a pullback. Hence the induced comparison
  \[
    (k^{[n+1]},h):U^{[n+1]}\to V^{[n+1]}
  \]
  is source-cartesian.
\end{enumerate}
\end{proof}

\begin{theorem}[Stability for finite paths]
\label{thm:stability-finite-paths-feedback}
Let
\[
  (k,h):U\to V
\]
be a source-cartesian comparison of endospans. For every
\[
  \psi\in\operatorname{Sub}(S'),
\]
one has
\[
  h^\ast
  \left(
    \langle\mathsf{Path}(V)\rangle\psi
  \right)
  =
  \left\langle
    \mathsf{Path}(U)
  \right\rangle
  h^\ast\psi,
\]
and
\[
  h^\ast
  \left(
    [\mathsf{Path}(V)]\psi
  \right)
  =
  \left[
    \mathsf{Path}(U)
  \right]
  h^\ast\psi.
\]
\end{theorem}

\begin{proof}
\begin{enumerate}
  \item \textbf{Prove the diamond equation for primitive lengths.}
  By Lemma~\ref{lem:iterated-source-cartesian-comparisons}, each finite
  iterate has a source-cartesian comparison
  \[
    (k^{[n]},h):U^{[n]}\to V^{[n]}.
  \]
  Lemma~\ref{lem:primitive-source-cartesian-modal-stability} gives
  \[
    h^\ast\langle V^{[n]}\rangle\psi
    =
    \langle U^{[n]}\rangle h^\ast\psi
  \]
  for every \(n\geq 0\).

  \item \textbf{Pass from finite lengths to paths for diamonds.}
  Since \(h^\ast\) preserves the countable joins used in the
  \(\omega\)-iterative semantics,
  \[
  \begin{aligned}
    h^\ast\langle\mathsf{Path}(V)\rangle\psi
    &=
    h^\ast
    \bigvee_{n\geq 0}
    \langle V^{[n]}\rangle\psi
    \\
    &=
    \bigvee_{n\geq 0}
    h^\ast\langle V^{[n]}\rangle\psi
    \\
    &=
    \bigvee_{n\geq 0}
    \langle U^{[n]}\rangle h^\ast\psi
    \\
    &=
    \langle\mathsf{Path}(U)\rangle h^\ast\psi.
  \end{aligned}
  \]

  \item \textbf{Repeat for boxes.}
  The box proof is dual. For every \(n\geq 0\),
  Lemma~\ref{lem:primitive-source-cartesian-modal-stability} gives
  \[
    h^\ast[ V^{[n]}]\psi
    =
    [U^{[n]}]h^\ast\psi.
  \]

  \item \textbf{Pass from finite lengths to paths for boxes.}
  Since \(h^\ast\) preserves the countable meets used in the
  \(\omega\)-iterative semantics,
  \[
  \begin{aligned}
    h^\ast[\mathsf{Path}(V)]\psi
    &=
    h^\ast
    \bigwedge_{n\geq 0}
    [V^{[n]}]\psi
    \\
    &=
    \bigwedge_{n\geq 0}
    h^\ast[V^{[n]}]\psi
    \\
    &=
    \bigwedge_{n\geq 0}
    [U^{[n]}]h^\ast\psi
    \\
    &=
    [\mathsf{Path}(U)]h^\ast\psi.
  \end{aligned}
  \]
\end{enumerate}
\end{proof}

\begin{corollary}[Stability for guarded while]
\label{cor:stability-guarded-while-feedback}
Let
\[
  (k,h):U\to V
\]
be a source-cartesian comparison of endospans
\[
  U:S\to S,
  \qquad
  V:S'\to S'.
\]
Let
\[
  \theta':\Theta'\hookrightarrow S',
  \qquad
  \epsilon':E'\hookrightarrow S'
\]
be guards, and let
\[
  \theta:=h^\ast\theta',
  \qquad
  \epsilon:=h^\ast\epsilon'
\]
be their pullbacks along \(h\). Then the induced comparison
\[
  U\circ T_\theta\longrightarrow V\circ T_{\theta'}
\]
is source-cartesian, and for every
\[
  \varphi\in\operatorname{Sub}(S')
\]
one has
\[
  h^\ast
  \left\langle
    (\theta'?;v)^{\star_{\mathrm{path}}};\epsilon'?
  \right\rangle\varphi
  =
  \left\langle
    (\theta?;u)^{\star_{\mathrm{path}}};\epsilon?
  \right\rangle
  h^\ast\varphi,
\]
and
\[
  h^\ast
  \left[
    (\theta'?;v)^{\star_{\mathrm{path}}};\epsilon'?
  \right]\varphi
  =
  \left[
    (\theta?;u)^{\star_{\mathrm{path}}};\epsilon?
  \right]
  h^\ast\varphi.
\]
Here \(u\) and \(v\) are program names whose interpretations are \(U\) and
\(V\), respectively.
\end{corollary}

\begin{proof}
\begin{enumerate}
  \item \textbf{Compare the guard tests.}
  Since
  \[
    \theta=h^\ast\theta',
    \qquad
    \epsilon=h^\ast\epsilon',
  \]
  the test spans \(T_\theta\) and \(T_\epsilon\) are obtained by pulling back
  the corresponding test spans \(T_{\theta'}\) and \(T_{\epsilon'}\) along
  \(h\).

  \item \textbf{Compare the guarded loop spans.}
  The comparison
  \[
    (k,h):U\to V
  \]
  and the pullback square defining \(T_\theta\) induce a source-cartesian
  comparison
  \[
    U\circ T_\theta
    \longrightarrow
    V\circ T_{\theta'}.
  \]
  This follows by composing source-cartesian comparisons and using pullback
  pasting.

  \item \textbf{Extend stability to finite paths.}
  Theorem~\ref{thm:stability-finite-paths-feedback} applies to the
  source-cartesian comparison of guarded loop spans and gives stability for
  \[
    \mathsf{Path}(U\circ T_\theta)
    \longrightarrow
    \mathsf{Path}(V\circ T_{\theta'}).
  \]

  \item \textbf{Compare the exit tests.}
  The exit tests are also related by pullback, so
  Lemma~\ref{lem:primitive-source-cartesian-modal-stability} gives the
  stability equation for the final \(T_\epsilon\)- and
  \(T_{\epsilon'}\)-steps.

  \item \textbf{Use the guarded while interpretations.}
  Using
  \[
    \llbracket
      (\theta?;u)^{\star_{\mathrm{path}}};\epsilon?
    \rrbracket
    =
    T_\epsilon\circ\mathsf{Path}(U\circ T_\theta)
  \]
  and
  \[
    \llbracket
      (\theta'?;v)^{\star_{\mathrm{path}}};\epsilon'?
    \rrbracket
    =
    T_{\epsilon'}\circ\mathsf{Path}(V\circ T_{\theta'}),
  \]
  the diamond equation follows by combining finite-path stability with
  Beck--Chevalley for the exit test.

  \item \textbf{Repeat for boxes.}
  The box equation is obtained by the same argument with universal
  quantification.
\end{enumerate}
\end{proof}

\section{Feedback after flattening response profiles}
\label{sec:feedback-after-flattening-response-profiles}

The polynomial response profile construction is one-step and coalgebraic.
Feedback is applied after flattening to an execution span.

Let
\[
  c:S\longrightarrow M_{A,B,Y}(S)
\]
be a relative response coalgebra, and write
\[
  \flat(c):S\to S
\]
for its flattened primitive execution endospan.

\begin{definition}[Profile trace]
\label{def:profile-trace-feedback}
The \textbf{profile trace} of \(c\) is
\[
  \operatorname{TrProfile}(c)
  :=
  \operatorname{Tr}^{S}(L_{\flat(c)}),
\]
where
\[
  L_{\flat(c)}
  =
  \begin{pmatrix}
    0_{S,S} & 1_S\\
    1_S & \flat(c)
  \end{pmatrix}.
\]
\end{definition}

\begin{proposition}[Profile trace after flattening]
\label{prop:profile-trace-after-flattening}
There is a canonical isomorphism
\[
  \operatorname{TrProfile}(c)
  \cong
  \mathsf{Path}(\flat(c)).
\]
\end{proposition}

\begin{proof}
\begin{enumerate}
  \item \textbf{Apply execution paths as feedback.}
  This is Theorem~\ref{thm:execution-paths-as-feedback} applied to the
  endospan
  \[
    \flat(c):S\to S.
  \]
\end{enumerate}
\end{proof}

For guards
\[
  \theta:\Theta\hookrightarrow S,
  \qquad
  \epsilon:E\hookrightarrow S,
\]
define
\[
  \operatorname{GTrProfile}_{\theta,\epsilon}(c)
  :=
  \operatorname{Tr}^{S}
  \left(
    \begin{pmatrix}
      0_{S,S} & 1_S\\
      T_\epsilon & \flat(c)\circ T_\theta
    \end{pmatrix}
  \right).
\]

\begin{proposition}[Guarded profile trace after flattening]
\label{prop:guarded-profile-trace-after-flattening}
There is a canonical isomorphism
\[
  \operatorname{GTrProfile}_{\theta,\epsilon}(c)
  \cong
  T_\epsilon\circ
  \mathsf{Path}(\flat(c)\circ T_\theta).
\]
Equivalently,
\[
  \operatorname{GTrProfile}_{\theta,\epsilon}(c)
  \cong
  \llbracket
    (\theta?;\flat(c))^{\star_{\mathrm{path}}};\epsilon?
  \rrbracket .
\]
\end{proposition}

\begin{proof}
\begin{enumerate}
  \item \textbf{Apply guarded while as guarded feedback.}
  This is Theorem~\ref{thm:guarded-while-as-guarded-trace} applied to the
  endospan
  \[
    \flat(c):S\to S.
  \]
\end{enumerate}
\end{proof}

Consequently,
\[
  \left\langle
    \operatorname{TrProfile}(c)
  \right\rangle\varphi
  =
  \mu X.
  \bigl(
    \varphi\vee\langle\flat(c)\rangle X
  \bigr),
\]
and
\[
  \left\langle
    \operatorname{GTrProfile}_{\theta,\epsilon}(c)
  \right\rangle\varphi
  =
  \mu X.
  \bigl(
    \epsilon\wedge\varphi
    \vee
    \theta\wedge\langle\flat(c)\rangle X
  \bigr).
\]
The corresponding box laws are
\[
  \left[
    \operatorname{TrProfile}(c)
  \right]\varphi
  =
  \nu X.
  \bigl(
    \varphi\wedge[\flat(c)]X
  \bigr),
\]
and
\[
  \left[
    \operatorname{GTrProfile}_{\theta,\epsilon}(c)
  \right]\varphi
  =
  \nu X.
  \bigl(
    (\epsilon\Rightarrow\varphi)
    \wedge
    (\theta\Rightarrow[\flat(c)]X)
  \bigr).
\]

\section{Examples}
\label{sec:feedback-trace-examples}

\subsection{Relations}
\label{subsec:feedback-relations}

When \(\mathcal E=\mathbf{Set}\), a span
\[
  X\xleftarrow{\ell}R\xrightarrow{r}Y
\]
has a relational image
\[
  x\,\overline R\,y
  \quad\Longleftrightarrow\quad
  \exists \rho\in R,\ \ell(\rho)=x,\ r(\rho)=y.
\]
Passing from spans to relational images sends homwise coproducts of spans to
unions of relations and span composition to relational composition.

Thus, for a relation
\[
  R\subseteq (X+H)\times(Y+H),
\]
the relational image of the categorical trace is
\[
  \overline{\operatorname{Tr}^{H}(R)}
  =
  \overline{R}_{XY}
  \cup
  \overline{R}_{HY}\circ
  \overline{R}_{HH}^{\ast}\circ
  \overline{R}_{XH},
\]
where \((-)^\ast\) is the ordinary relational Kleene star, not pullback and not
the syntactic operation \((-)^{\star_{\mathrm{path}}}\). The bar notation
emphasizes that the span trace itself retains witnesses, while the relational
image forgets them.

\subsection{Classical Mealy machines}
\label{subsec:feedback-classical-mealy}

For an ordinary Mealy machine with state set \(S\), the primitive execution
span
\[
  U:S\to S
\]
records one-step executions. The finite-path endospan
\[
  \mathsf{Path}(U)
\]
records finite executions of arbitrary length. By
Theorem~\ref{thm:execution-paths-as-feedback},
\[
  \mathsf{Path}(U)
  \cong
  \operatorname{Tr}^{S}
  \begin{pmatrix}
    0_{S,S} & 1_S\\
    1_S & U
  \end{pmatrix}.
\]
Thus finite execution is a coproduct feedback trace.

\subsection{Boolean guarded while}
\label{subsec:boolean-guarded-while}

In a Boolean subobject fibre, take the exit predicate to be
\[
  \epsilon=\neg\theta.
\]
Then the dynamic program
\[
  (\theta?;u)^{\star_{\mathrm{path}}};(\neg\theta)?
\]
has diamond semantics
\[
  \left\langle
    (\theta?;u)^{\star_{\mathrm{path}}};(\neg\theta)?
  \right\rangle\varphi
  =
  \mu X.
  \bigl(
    \neg\theta\wedge\varphi
    \vee
    \theta\wedge\langle u\rangle X
  \bigr).
\]
The box semantics is
\[
  \left[
    (\theta?;u)^{\star_{\mathrm{path}}};(\neg\theta)?
  \right]\varphi
  =
  \nu X.
  \bigl(
    (\neg\theta\Rightarrow\varphi)
    \wedge
    (\theta\Rightarrow[u]X)
  \bigr).
\]

\subsection{Signed interfaces and Geometry of Interaction}
\label{subsec:signed-interfaces-feedback}

An Int morphism
\[
  (X^+,X^-)\longrightarrow(Y^+,Y^-)
\]
is a span
\[
  X^+ + Y^-
  \longrightarrow
  Y^+ + X^-.
\]
Thus the positive input of the source and the negative input of the target are
read together, while the positive output of the target and the negative output
of the source are produced together. Composition feeds the middle signed
interface
\[
  Y^+ + Y^-
\]
back into itself and hides it by trace. This is the compact-closed form of
feedback span semantics.

This interface discipline is the same structural pattern that appears in
Geometry of Interaction: computation is represented by interaction through an
internal boundary, and composition hides the boundary by an execution
operation
\cite{Girard1989GeometryInteractionI,AbramskyHaghverdiScott2002GeometryInteractionLCA}.
In the present chapter the execution operation is the finite hidden-path trace
in a coproduct-monoidal span category.

\subsection{Terminal downstream response profiles}
\label{subsec:terminal-response-profiles-feedback}

In the terminal downstream case, flattening a terminal response coalgebra gives
the terminal primitive execution span. The trace construction gives the
finite-path endospan of that span:
\[
  \operatorname{TrProfile}(c)
  \cong
  \mathsf{Path}(\flat(c)).
\]
Thus the terminal downstream profile trace is the one-step profile
presentation of the same finite feedback semantics after flattening.

\section{Chapter summary}
\label{sec:feedback-trace-summary}

This chapter constructed feedback as the trace for the coproduct monoidal
structure on spans.

The foundational structure is the coproduct symmetric monoidal bicategory
\[
  \bigl(
    \mathbf{Span}^{+}_{\omega}(\mathcal E),+,0
  \bigr).
\]
Passing to isomorphism classes gives the symmetric monoidal category
\[
  \bigl(
    \mathbf{Span}^{+}_{\omega,0}(\mathcal E),+,0
  \bigr).
\]

For an endospan
\[
  U:S\to S,
\]
the finite-path endospan is
\[
  \mathsf{Path}(U)
  =
  \coprod_{n\geq 0}U^{[n]}.
\]
It is the free endospan monoid on \(U\), and the assignment
\[
  U\mapsto\mathsf{Path}(U)
\]
is the path monad on endospans. The monoid unit is the inclusion of the
\(0\)-summand
\[
  1_S\to\mathsf{Path}(U),
\]
while the monad unit is the inclusion of the \(1\)-summand
\[
  U\to\mathsf{Path}(U).
\]

For a span
\[
  R:X+H\to Y+H
\]
with block decomposition
\[
  R=
  \begin{pmatrix}
    R_{XY} & R_{XH}\\
    R_{HY} & R_{HH}
  \end{pmatrix},
\]
the coproduct feedback trace is
\[
  \operatorname{Tr}^{H}(R)
  =
  R_{XY}
  \boxplus
  R_{HY}\circ
  \mathsf{Path}(R_{HH})
  \circ
  R_{XH}.
\]
This makes
\[
  \bigl(
    \mathbf{Span}^{+}_{\omega,0}(\mathcal E),+,0
  \bigr)
\]
a traced symmetric monoidal category, and gives hom-category trace functors at
the bicategorical level.

The trace induces a Conway dagger. For
\[
  f:X\to Y+X
\]
with blocks
\[
  f_Y:X\to Y,
  \qquad
  f_X:X\to X,
\]
one has
\[
  f^\dagger
  =
  f_Y\circ\mathsf{Path}(f_X).
\]
This dagger satisfies the Conway fixpoint law
\[
  f^\dagger
  \cong
  f_Y\boxplus(f^\dagger\circ f_X),
\]
as well as output naturality, uniformity, sliding, and the double-dagger
identity.

The same coproduct tensor lifts to internal categories, Mealy morphisms, and
Mealy simulations, giving a coproduct symmetric monoidal bicategory
\[
  \mathbf{Mly}^{+}_{\omega}(\mathcal E).
\]
The trace construction lifts to Mealy morphisms:
\[
  \operatorname{Tr}^{C}(\Phi)
  =
  \Phi_{AB}
  \boxplus
  \Phi_{CB}\circ
  \mathsf{Path}(\Phi_{CC})
  \circ
  \Phi_{AC}.
\]
It also acts on simulations by taking finite hidden-path composites of the
corresponding block simulations:
\[
  \operatorname{Tr}^{C}(\Theta)
  =
  \Theta_{AB}
  \boxplus
  \coprod_{n\geq 0}
  \left(
    \Theta_{CB}
    \ast_{\mathrm h}
    \Theta_{CC}^{[n]}
    \ast_{\mathrm h}
    \Theta_{AC}
  \right),
\]
where the order follows the standing categorical convention: first
\(\Theta_{AC}\), then \(\Theta_{CC}^{[n]}\), then \(\Theta_{CB}\).

The Int construction packages feedback into compact closed signed-interface
categories. For spans, an Int morphism
\[
  (X^+,X^-)\to(Y^+,Y^-)
\]
is a span
\[
  X^+ + Y^-\to Y^+ + X^-.
\]
For Mealy morphisms, an Int morphism
\[
  (A^+,A^-)\to(B^+,B^-)
\]
is a Mealy morphism
\[
  A^+ + B^-
  \rightsquigarrow
  B^+ + A^-.
\]
Composition traces over the middle signed interface.

Relative execution paths are feedback traces. For a sub-endospan
\[
  U\hookrightarrow m_{\Phi,Y},
\]
one has
\[
  \operatorname{Tr}^{S}
  \begin{pmatrix}
    0_{S,S} & 1_S\\
    1_S & U
  \end{pmatrix}
  \cong
  \mathsf{Path}(U).
\]
The folded execution map is the unique monoid morphism
\[
  \mathsf{Path}(U)\to m_{\Phi,Y}
\]
induced by the inclusion \(U\hookrightarrow m_{\Phi,Y}\).

In the logical setting, finite-path modalities satisfy
\[
  \langle\mathsf{Path}(U)\rangle\varphi
  =
  \mu X.
  \bigl(
    \varphi\vee\langle U\rangle X
  \bigr),
\]
and
\[
  [\mathsf{Path}(U)]\varphi
  =
  \nu X.
  \bigl(
    \varphi\wedge[U]X
  \bigr).
\]
Consequently, feedback trace has diamond semantics
\[
  \left\langle
    \operatorname{Tr}^{H}(R)
  \right\rangle\varphi
  =
  \langle R_{XY}\rangle\varphi
  \vee
  \langle R_{XH}\rangle
  \mu Z.
  \bigl(
    \langle R_{HY}\rangle\varphi
    \vee
    \langle R_{HH}\rangle Z
  \bigr),
\]
with the dual greatest-fixed-point box law.

Dynamic iteration is interpreted by
\[
  \llbracket p^{\star_{\mathrm{path}}}\rrbracket
  =
  \mathsf{Path}(\llbracket p\rrbracket),
\]
and tests are interpreted by
\[
  \llbracket\theta?\rrbracket=T_\theta.
\]
For guarded while programs,
\[
  \llbracket
    (\theta?;u)^{\star_{\mathrm{path}}};\epsilon?
  \rrbracket
  =
  T_\epsilon\circ
  \mathsf{Path}(\llbracket u\rrbracket\circ T_\theta).
\]
Hence
\[
  \left\langle
    (\theta?;u)^{\star_{\mathrm{path}}};\epsilon?
  \right\rangle\varphi
  =
  \mu X.
  \bigl(
    \epsilon\wedge\varphi
    \vee
    \theta\wedge\langle u\rangle X
  \bigr),
\]
and
\[
  \left[
    (\theta?;u)^{\star_{\mathrm{path}}};\epsilon?
  \right]\varphi
  =
  \nu X.
  \bigl(
    (\epsilon\Rightarrow\varphi)
    \wedge
    (\theta\Rightarrow[u]X)
  \bigr).
\]

The stability theorem uses source-cartesian comparisons of execution
endospans. If
\[
  (k,h):U\to V
\]
is source-cartesian, then finite-path modalities commute with pullback along
\(h\):
\[
  h^\ast
  \left(
    \langle\mathsf{Path}(V)\rangle\psi
  \right)
  =
  \left\langle
    \mathsf{Path}(U)
  \right\rangle
  h^\ast\psi,
\]
and
\[
  h^\ast
  \left(
    [\mathsf{Path}(V)]\psi
  \right)
  =
  \left[
    \mathsf{Path}(U)
  \right]
  h^\ast\psi.
\]
This gives stability for guarded while programs under downstream change and
Mealy simulation when the primitive execution spans and guards are transported
by the corresponding source-cartesian comparison.

Finally, polynomial response profiles enter feedback only after flattening:
\[
  \operatorname{TrProfile}(c)
  \cong
  \mathsf{Path}(\flat(c)),
\]
and
\[
  \operatorname{GTrProfile}_{\theta,\epsilon}(c)
  \cong
  T_\epsilon\circ
  \mathsf{Path}(\flat(c)\circ T_\theta).
\]
Thus coproduct feedback trace, finite execution paths, Conway iteration,
guarded while programs, signed-interface Int semantics, Geometry of
Interaction, and fixed-point dynamic semantics are all presentations of the
same finite hidden-path construction.